\documentclass[usenatbib,twoside,twocolumn]{mnras}

% Paper packages
\usepackage{graphicx}
\usepackage{amsmath}
\usepackage{amssymb}
\usepackage{booktabs}

% Denser float packing (reduce float-only pages; no content change)
\renewcommand{\topfraction}{0.9}
\renewcommand{\bottomfraction}{0.85}
\renewcommand{\textfraction}{0.07}
\renewcommand{\floatpagefraction}{0.75}
\renewcommand{\dbltopfraction}{0.9}
\renewcommand{\dblfloatpagefraction}{0.75}
\setcounter{topnumber}{3}
\setcounter{bottomnumber}{2}
\setcounter{totalnumber}{5}
\setcounter{dbltopnumber}{3}

% Define astronomical units
\newcommand{\msun}{\ensuremath{M_{\odot}}}
\newcommand{\pc}{\ensuremath{\rm pc}}
\newcommand{\kms}{\ensuremath{\rm km\,s^{-1}}}

\title[Sub-classical core spacing in HGBS filaments]{Core Spacing in Herschel Gould Belt Survey Filaments: a Convention-Limited Sub-Classical Scale and the Limits of Its Interpretation}

\author[G.~J.~White]{G.~J.~White$^{1,2}$\\
$^{1}$School of Physical Sciences, The Open University, Walton Hall, Milton Keynes, MK7 6AA, UK\\
$^{2}$Rutherford Appleton Laboratory, Harwell Science and Innovation Campus, Didcot, OX11 0QX, UK}

\begin{document}

\maketitle

\begin{abstract}
\textit{Herschel} finds dense cores spaced along filaments well below the classical Inutsuka \& Miyama fragmentation wavelength ($\lambda_{\rm m}\approx5$--$6\times$FWHM). We re-analyse core spacings in \emph{four} nearby, low-mass Herschel Gould Belt Survey clouds with Gaia DR3 distances. With only $N=4$ clouds our numbers are a descriptive characterisation of this regime, not a population estimate, and we make no claim beyond nearby low-mass star-forming gas. Our robust findings are methodological: the nearest-neighbour (NN) spacing is the appropriate observable, whereas the widely-used pairwise median is biased, converging to the region extent $L/3$; and the NN spacing is sub-classical in \emph{sign} under every skeleton and width convention. Its magnitude is convention-dominated --- the skeleton algorithm alone (DisPerSE vs FilFinder) moves $\lambda/W$ by a factor of two, the single largest uncertainty --- so the defensible physical statement is $\lambda/W\ll$ classical, not a specific value (we quote $(\lambda/W)_0\approx1.4$ as a reference only). The cores reject \emph{both} periodic and pure-Poisson spacing (parametric-bootstrap calibrated), consistent with a hard-core-excluded line density with region-dependent clustering, rather than a resonant wavelength. A large Athena++ MHD campaign shows a dynamically contracting filament \emph{can} fragment sub-classically --- a resolution-stable upper bound, validated by an equilibrium-recovery control that recovers the classical wavelength; this demonstrates plausibility, not a unique cause. We do not identify the mechanism: single-filament dynamical fragmentation, hierarchical fibres, and (in the untested force-balanced helical geometry) magnetic tension all remain open. The contribution is primarily methodological.
\end{abstract}


\begin{keywords}
stars: formation -- ISM: structure -- ISM: filaments -- MHD -- turbulence
\end{keywords}

\section{INTRODUCTION}

The \textit{Herschel} Space Observatory established that filaments are the fundamental structures within which star formation occurs in molecular clouds \citep{Andre2010}. Two key observational results emerged: (1) a characteristic filament width of $\sim$0.1 pc across diverse environments \citep{Arzoumanian2011}, and (2) a critical mass-per-unit-length of $M_{\rm line,crit} \approx 16~M_\odot$/pc \citep{Ostriker1964}.

Classical fragmentation theory \citep{Inutsuka1992} predicts that isothermal gas cylinders fragment at a wavelength $\lambda_{\rm m}\approx5$--$6\times$ the filament FWHM (\citealt{Fischera2012}; equivalently $4\times$ the flat inner diameter, \citealt{Inutsuka1992}); for the characteristic HGBS width of $0.10$\,pc this is a spacing of $\sim0.5$\,pc. Our adopted nearest-neighbour spacing is $(\lambda/W)_0\approx1.4$, using directly-fitted region-specific Plummer widths (Section~\ref{sec:widths}); the fixed-width convention gives $\lambda/W\approx2.1$ (raw; $1.9$ after the periodic injection--recovery correction), retained only for comparability with prior HGBS work (Table~\ref{tab:headline} decodes all correction states; both are sub-classical by a factor $\approx3$--$4$). We adopt Gaia DR3 distances \citep{Zhang2023} and the nearest-neighbour (NN) statistic, the estimator that most directly measures the fragmentation wavelength and the one used in the HGBS catalogue papers \citep{Konyves2015,Konyves2020} (Section~\ref{sec:statistics}).

This apparent deficit relative to the classical $\lambda_{\rm m}\approx5$--$6\times$FWHM has been attributed to additional physics. \citet{Arzoumanian2019} established the $\sim$0.1\,pc \emph{width}, not a spacing; the sub-classical \emph{spacing} deficit and its usual fibre/magnetic reading come from the fragmentation literature (e.g.\ \citealt{Fischera2012,Zhang2023cal}). Two mechanisms are usually invoked: \emph{hierarchical fragmentation}, in which filaments are bundles of velocity-coherent fibres \citep{Hacar2013,Hacar2018,Yang2024} so that measuring core spacings across the bundle compresses the apparent filament-level spacing; and \emph{magnetic tension} \citep{Nakamura1993}. Only a \emph{helical}/toroidal field shortens the axial mode --- a static longitudinal field leaves it unchanged and a perpendicular field lengthens it (Section~\ref{sec:magnetic}) --- so ``magnetic tension'' as a route to a shorter spacing is a specifically helical effect whose force-balanced case remains untested (Section~\ref{sec:optional_tests}). These are natural readings within equilibrium fragmentation theory; we suggest instead that the appropriate comparison for a still-collapsing filament is time-dependent, which reduces the deficit without invoking them, though it does not by itself select among the readings.

From a large campaign of self-gravitating MHD simulations we measure the fragmentation wavelength of a dynamically collapsing filament from its growth-rate spectrum, subject to two caveats. First, the measured wavelength is an order-$\lambda_{\rm J}$ \emph{upper bound}, not a located peak (throughout, $\lambda_{\rm J}=c_s(\pi/G\rho_0)^{1/2}$ is the thermal Jeans length at the background density $\rho_0$, so that the classical cylinder wavelength is $\lambda_{\rm m}=22H=2.3\,\lambda_{\rm J}$; $H=c_s/\sqrt{4\pi G\rho_c}$): the spectrum rises to a plateau at the short-wavelength edge of the resolution-converged band ($\lambda\gtrsim1\,\lambda_{\rm J}$) and does not turn over, so $\lambda_{\rm max}\lesssim1\,\lambda_{\rm J}$ and the true dynamical scale may be shorter (Appendix~\ref{sec:val_appendix}). Second, the like-for-like comparison must use matched widths: in a common beam-convolved Plummer convention the simulated value is $\lambda/W\approx1.0$--$1.3$, between the FilFinder ($0.4$--$0.8$) and DisPerSE ($1.9$, periodic-corrected; Table~\ref{tab:headline}) observed values. These simulations therefore demonstrate that a contracting filament \emph{can} fragment sub-classically \citep{Clarke2016,Clarke2017}; they establish plausibility and an upper bound, not that dynamical fragmentation \emph{explains} the observed spacing. Because our seeded profile is not the exact isothermal equilibrium (Section~\ref{sec:limitations}), we cannot fully separate the contraction effect from an initial-condition contribution. The observations are \emph{consistent with} single-filament dynamical fragmentation but do not \emph{require} it, nor discriminate it from hierarchical fibres or magnetic tension.

We present a full-coverage nearest-neighbour core-spacing analysis with Gaia DR3 distances \citep{Zhang2023} for the four robust HGBS regions (Orion~B, Aquila, Perseus, Taurus; the eight-region sample is retained only for the legacy pairwise-median statistic), followed by the simulation campaign, its validation and the interpretation. We note throughout that the large-scale filament--field misalignment from Planck polarimetry \citep{Planck2016} measures the projected orientation, which need not equal a filament's internal field geometry (Section~\ref{sec:discussion}).

\smallskip\noindent\emph{What is and is not established.} Two groups of statements should be kept apart. \emph{Established here:} the pairwise median is a biased estimator that converges to the region extent $L/3$; the nearest-neighbour spacing is the appropriate observable; the spacing is sub-classical in \emph{sign} under every skeleton and width convention; and its magnitude is convention-dominated, the skeleton-algorithm choice (DisPerSE vs FilFinder) being the single largest term. \emph{Not established:} that dynamical collapse is the dominant cause; that hierarchical fibres are unimportant; and that magnetic tension is excluded. In notation used throughout (Table~\ref{tab:symbols}), $\lambda$ is the observed core-to-core spacing, $\lambda_{\rm max}$ the simulated dominant fragmentation wavelength, $\lambda_{\rm m}$ the classical equilibrium wavelength, $W\equiv W_{\rm fil}$ the observed filament FWHM, and $W_{\rm core}$ the simulation initial half-width.

\begin{table}
\caption{Principal symbols.}
\label{tab:symbols}
\footnotesize
\setlength{\tabcolsep}{4pt}
\begin{tabular}{lp{6.1cm}}
\toprule
Symbol & Meaning \\
\midrule
$\lambda$ & core-to-core spacing (an observed line-density scale, not necessarily a resonant wavelength; cf.\ Section~\ref{sec:clustering}) \\
$W_{\rm fil}$ & observed filament FWHM (beam-deconvolved Plummer; the reference width) \\
$W_{\rm core}$ & simulation initial Gaussian half-width $\sigma$ ($0.3\,\lambda_{\rm J}$) \\
$W_{\rm form}$ & FWHM of the evolved, projected simulation profile \\
$T_1\!\equiv\!\eta_W$ & $W_{\rm form}/W_{\rm fil}$, the simulation-to-observation width normalisation (forward-model band $0.59$--$0.78$; adopted $0.61$) \\
$\lambda_{\rm J}$ & background Jeans length ($=1$ in code units) \\
$\lambda_{\rm max}$ & dominant fragmentation wavelength (growth-rate peak) \\
$\lambda_{\rm IM92}$ & classical equilibrium-cylinder wavelength $\approx5$--$6\,W_{\rm FWHM}$ ($=4\times$ flat diameter) \\
$\Gamma(\lambda)$ & linear growth rate of axial mode of wavelength $\lambda$ \\
$f$ & thermal line-mass fraction $M_{\rm line}/M_{\rm line,crit}$ \\
$\beta$ & plasma beta $8\pi P/B^2$ \\
$\mathcal{M}$ & amplitude label of the seed perturbation ($\delta v=\mathcal{M}c_s\times10^{-4}$; a linear seed, \emph{not} a turbulent Mach number --- driven transonic turbulence is tested separately in Appendix~\ref{sec:supplementary_campaigns}) \\
$\mu_\Phi/\mu_{\Phi,\rm crit}$ & normalised mass-to-flux ratio \\
$\theta$ & angle between field and filament axis \\
$N_J$ & cells per local Jeans length (Truelove criterion) \\
$C$ & density contrast $\rho_{\rm max}/\langle\rho\rangle$ \\
\bottomrule
\end{tabular}
\end{table}

\section{OBSERVATIONAL ANALYSIS}

\subsection{Complete HGBS Sample}

Our sample comprises 8 high-confidence HGBS regions (Table~\ref{tab:sample}). The original catalogues \citep{Konyves2015, Konyves2020, Ladjelate2020, Pezzuto2021} adopted distances of $130$--$400$\,pc from the best estimates available in 2010--2015 (Table~\ref{tab:literature} lists the pre-Gaia distance and its source per region; note that \citealt{Konyves2020} already adopted $\sim400$\,pc for Orion~B); Gaia DR3 parallaxes \citep{Zhang2023} have since enabled substantial revisions, largest for the more distant regions, whose pre-Gaia distances were the most uncertain (Gaia fractional parallax errors are in fact largest, not smallest, at greater distance).

We adopt the Gaia DR3 YSO-based distances from \citet{Zhang2023} (Table~4), with typical uncertainties of 3\%--5\%. The revisions are strongly region-dependent: largest for Serpens ($260\to458$\,pc, $+76\%$) and Aquila ($260\to436$\,pc, $+68\%$), whereas Orion~B, the largest region, is essentially unchanged (pre-Gaia $\sim400$\,pc, \citealt{Konyves2020}; Gaia DR3 $386$\,pc, $-3.5\%$), and Taurus ($140\to135$\,pc) and Ophiuchus ($130\to137$\,pc) barely move. Since spacings scale linearly with distance, the core-count-weighted mean rises only modestly ($\sim10\%$ for the robust regions), dominated by Aquila rather than being a uniform $+32\%$ shift. The observed spacing is a line-density statistic, not a resonance, and coincides with the theoretical wavelength only to order of magnitude.

The analysis includes all HGBS-identified cores (starless, prestellar and protostellar), extracted and classified within the consistent HGBS framework \citep{Konyves2015, Konyves2020}. Including unbound starless and protostellar sources may introduce some bias, but fragmentation produces a continuum of boundness states, and restricting the sample to prestellar cores would shrink several regions (notably TMC1, CRA and Serpens) below statistical usefulness. Core positions and properties are taken directly from the published catalogues.

Protostellar migration displacements of $0.01$--$0.05$\,pc \citep{Kirk2016, Mattern2018} are small against the $\sim0.2$--$0.3$\,pc spacing; we estimate the resulting NN bias at $\sim5$--$10\%$ and carry it as a systematic (Section~\ref{sec:error_budget}).

\begin{table*}
\caption{Complete HGBS Sample with Gaia DR3 Distances. Per-region spacings are pairwise-median values rounded to three significant figures; the core-count-weighted mean is $0.279$\,pc (full sample) and $0.290$\,pc (robust), the latter used only for the robust-vs-full comparison (Section~2.3).}
\label{tab:sample}
\begin{tabular}{lccccc}
\toprule
Region & Distance & Total & Pairwise median & Std Error & Status \\
       & (pc)   & Cores & (pc)   & (pc)      &  \\
\midrule
Orion B  & 386 & 1,844 & 0.313 & 0.047 & Robust \\
Aquila   & 436 & 749   & 0.346 & 0.047 & Robust \\
Perseus  & 300 & 816   & 0.248 & 0.040 & Robust \\
Taurus   & 135 & 536   & 0.198 & 0.040 & Robust \\
Ophiuchus& 137 & 513   & 0.206 & 0.053 & Limited \\
Serpens  & 458 & 194   & 0.331 & 0.097 & Limited \\
TMC1     & 135 & 178   & 0.195 & 0.056 & Limited \\
CRA      & 150 & 239   & 0.248 & 0.072 & Limited \\
\midrule
\textbf{Weighted Mean}$^\dagger$ &  & \textbf{5,069} & \textbf{0.279} & \textbf{0.009} &  \\
\bottomrule
\end{tabular}
\par\vspace{2pt}
\footnotesize $^\dagger$The $0.009$\,pc figure is a formal core-count-weighted standard error (treating cores as independent); it is \emph{not} the relevant uncertainty, since the four clouds are the independent units (Section~\ref{sec:error_budget}). It is retained only as a within-region formal quantity and plays no role in the comparison with theory.
\end{table*}

\subsection{Measurement Methodology}
\label{sec:methodology}

The analysis pipeline, from raw HGBS maps to the final characteristic $\lambda/W$, is summarised in Fig.~\ref{fig:pipeline}. Filament skeletons were identified using DisPerSE \citep{Sousbie2011} following standard HGBS methodology \citep{Arzoumanian2011}, with a $3\sigma$ persistence threshold and manual editing restricted to removing obvious artefacts (e.g.\ spurs connecting unrelated structures). Its effect on the spacing is small ($\lesssim5\%$), bounded using the junction-retention and association-threshold sensitivity tests as a \emph{proxy} (Table~\ref{tab:errorbudget}); a blind, multi-editor re-skeletonisation is the direct confirmation and remains outstanding (Section~\ref{sec:limitations}). Because skeletons and core positions derive from distance-independent angular positions, the Gaia DR3 revisions require no re-extraction. Cores at filament junctions are retained (standard HGBS practice) and associated with filaments within a perpendicular threshold of $2\times$ the characteristic width ($0.2$\,pc). The measured spacing is insensitive to these choices: excluding junctions changes it by $<5\%$ and varying the association threshold between $1\times$ and $1.5\times$ the width by $<3\%$. (This $2\times$-width threshold refers to the pairwise-median association; the primary along-skeleton NN pipeline associates cores within a half-width, $\sim0.06$\,pc, of a spine, Section~\ref{sec:injection_recovery}.)

Core-to-core spacings were also measured using the pairwise median statistic, which Section~\ref{sec:statistics} shows converges to $L/3$ rather than the fragmentation wavelength for $N\gtrsim5$; we retain it only as a comparability statistic with prior HGBS work.

Because physical spacing is angular separation times distance, the $\lambda\propto d$ scaling is tautological and we do not present it as corroboration; the relevant check is that the \emph{residuals} about that scaling do not correlate with the absolute revision fraction $|\Delta D/D_{\rm original}|$ ($r=0.18$, $p=0.68$), so the revision does not preferentially bias the most-revised clouds. Two caveats remain, and we flag them as limitations rather than resolve them: the YSO-based distances \citep{Zhang2023} are assumed to apply to the dense gas measured here, and the Aquila/Serpens Rift is known to have substantial line-of-sight depth, so assigning a single distance may blend populations and inflate the inferred spacing. As a floor, retaining the pre-Gaia distances lowers the robust-region characteristic to $\lambda/W\approx1.4$ (still sub-classical), so the sub-classical \emph{sign} does not depend on the distance revision, though its magnitude does.

\begin{figure}
\centering
\begin{minipage}{0.95\columnwidth}
\footnotesize
\begin{enumerate}\setlength{\itemsep}{2pt}
\item HGBS column-density maps ($250\,\mu$m reference) + Gaia DR3 distances
\item DisPerSE skeleton extraction ($3\sigma$ persistence threshold)
\item Morphological closing (8 iterations) $+$ re-skeletonisation (bridge fragmented segments; components $\geq40$\,px)
\item Core association: cores within $0.06$\,pc of a spine are assigned; others rejected
\item Graph ordering: cores ordered by double-BFS arclength along each component
\item Adjacent-core NN spacing (in pc, via Gaia DR3 distance)
\item Injection--recovery bias correction ($\div\,1.11$--$1.29$, region-dependent)
\item Characteristic $\lambda/W = $ bias-corrected median spacing $/\,W_{\rm fil}$
\end{enumerate}
\end{minipage}
\caption{Analysis pipeline from raw HGBS maps to the final characteristic $\lambda/W$. Each step's parameters are given in the text (Sections~\ref{sec:statistics}--\ref{sec:injection_recovery}).}
\label{fig:pipeline}
\end{figure}

\subsection{Results}

The four robust regions are the independent statistical units throughout this paper; individual cores and spacings within a cloud are not independent, and we propagate this nesting in the hierarchical bootstrap (Section~\ref{sec:error_budget}). With only four independent clouds we treat the cloud-to-cloud scatter ($\sigma\approx0.4$) as a descriptive range, not a formal confidence interval; we do not attach a bootstrap CI to an $N=4$ sample. As an illustration of how weakly four points constrain a population mean --- not a claimed confidence interval --- a Student-$t$ interval on the mean with three degrees of freedom would span mean\,$\pm\,t_{3,0.975}\,\sigma/\sqrt{4}$ ($t_{3,0.975}=3.18$), i.e.\ $\pm\sim0.6$, comparable to the central value itself.

\begin{table}
\caption{Convention decoder: read every $\lambda/W$ quoted in the text against it. The nearest-neighbour $\lambda/W$ under each width/skeleton convention and correction state, for the four robust HGBS regions (revised distances); every combination is sub-classical (classical $\lambda_{\rm m}\approx5$--$6\times$FWHM). We adopt the directly-measured region-Plummer value as our reference, $\boldsymbol{(\lambda/W)_0\approx1.4}$; the fixed-width value ($\approx1.9$, retained for comparability with prior HGBS work) and all others are read against it. The ``periodic'' column applies the periodic injection--recovery correction to the raw value; the data-matched \emph{clustered} column (bracketed) applies $\div1.64$ to the \emph{raw} value, giving $\approx1.2$ (fixed width) and $\approx0.8$ (region-Plummer). The DisPerSE conventions are sub-classical by a factor $\approx3$--$4$ (FilFinder more).}
\label{tab:headline}
\footnotesize
\setlength{\tabcolsep}{4pt}
\begin{tabular}{p{3.6cm}ccc}
\toprule
Convention & raw & periodic & clustered$^{\dagger}$ \\
\midrule
Region Plummer, DisPerSE (adopted, $(\lambda/W)_0$) & $\mathbf{1.4}$ & $1.2$ & $[0.8]$ \\
Fixed $0.10$\,pc, DisPerSE & $2.1$ & $1.9$ & $[1.2]$ \\
FilFinder skeleton & $0.5$--$0.9$ & $0.4$--$0.8$ & $[0.3$--$0.5]$ \\
Pairwise median ($L/3$, not $\lambda$) & $2.8$ & --- & --- \\
\midrule
Matched simulation (upper limit) & \multicolumn{3}{c}{$\lambda/W\lesssim1.0$--$1.3$} \\
Classical $\lambda_{\rm m}$ (F\&M 2012) & \multicolumn{3}{c}{$5$--$6$ ($=4\times$ diam.)} \\
\bottomrule
\end{tabular}
\\[2pt]
\footnotesize
\emph{Notes} The adopted value is the region-Plummer NN $(\lambda/W)_0\approx1.4$ (bold; plausible range $0.8$--$1.9$ across the conventions and injection--recovery corrections; not a formal confidence interval; Section~\ref{sec:error_budget}). $^{\dagger}$The \emph{clustered} column is not an independent estimate: it applies a correction derived by \emph{assuming} the clustering the data are observed to have (Section~\ref{sec:injection_recovery_model}), so it is circular and shown only as a lower bound. We carry the raw value throughout. The fixed-width convention gives a \emph{higher} value ($\approx1.9$) because it uses $0.10$\,pc rather than the larger measured widths; FilFinder gives lower values still. The pairwise median converges to $0.293\,L\approx0.88\,(L/3)$ (Section~\ref{sec:statistics}) --- a region-extent scale --- overestimating the fragmentation wavelength by $\sim$45--50\%. We quote $\lambda/W$ to one decimal place; finer digits are not significant.
\end{table}

\textbf{Robust regions.} We adopt the four ``robust'' regions (Orion~B, Aquila, Perseus, Taurus; $N>500$ cores each, well-established distances) as the primary sample.

\textbf{Full sample result.} The core-count-weighted mean across all 8 HGBS regions ($0.279$\,pc) and the robust-only weighted mean ($0.290$\,pc) differ by $\approx3.9\%$, and a leave-one-out analysis confirms no single region dominates ($<6\%$ shifts). All spacings are quoted to three significant figures; where an exact recomputation from the rounded per-region inputs differs from the bootstrap-derived primary value it does so by $\lesssim1\%$, which does not affect any conclusion.

All weighted means use core-count weighting,
\begin{equation}
\bar{\lambda} = \frac{\sum_i N_i\,\lambda_i}{\sum_i N_i},
\label{eq:weighting}
\end{equation}
where $\lambda_i$ and $N_i$ are the per-region spacing and core count.

With only four independent clouds we do not attach a formal confidence level to the comparison: a nonparametric resampling of four regions cannot estimate a population variance, so we quote the four per-region values and their min--max range rather than a bootstrap ``95\% CI'' (resampling the four pairwise-median values spans $0.261$--$0.298$\,pc, descriptive only). The absolute $\lambda/W$ result below uses the adopted Gaia distances. The preferred NN measurement is presented in Section~\ref{sec:injection_recovery}.

\subsection{The dominant observational systematic: filament width}
\label{sec:widths}

The fixed $W_{\rm fil}=0.10$\,pc convention is the largest single lever on $\lambda/W$, comparable to the whole discrepancy, so we measured region-specific widths directly. Along each spine we sampled the column-density profile perpendicular to the local tangent and fitted a beam-deconvolved Plummer function (the HGBS width definition; \citealt{Arzoumanian2019}) to hundreds of profiles per region ($18.2''$ beam; Table~\ref{tab:widths}). All four regions give $p\simeq2.0$ and $W_{\rm fil}=0.11$--$0.16$\,pc, agreeing with the published HGBS values and slightly broader than $0.10$\,pc, so the region-specific $\lambda/W$ ($1.1$--$1.8$, median $\approx1.4$) is \emph{lower} than the raw fixed-width value ($\approx2.1$; Table~\ref{tab:headline}); every region is sub-classical under either convention. A Gaussian-crest estimate returns a spuriously narrow Taurus width ($0.031$\,pc, formally $\lambda/W\to4$); the physically-motivated Plummer fit ($0.11$\,pc, $\lambda/W\approx1.1$) makes Taurus the most sub-classical region. The result is thus not conditional on the width convention.

\begin{table}
\caption{Region-specific, beam-deconvolved filament widths from a direct Plummer fit to perpendicular column-density profiles along the skeletons (the HGBS width definition; \citealt{Arzoumanian2019}), and the resulting NN $\lambda/W$. $W_{\rm fil}$ is the beam-deconvolved Plummer FWHM ($18.2''$ beam), $p$ the Plummer index (a \emph{free} fit parameter, not fixed; it converges to $p=2.00$--$2.01$ in all four regions, consistent with the canonical Plummer-2 profile), $N$ the number of fitted profiles. All four regions give $p\simeq2$ and sub-classical $\lambda/W$.}
\label{tab:widths}
\footnotesize
\setlength{\tabcolsep}{3pt}
\begin{tabular}{p{1.7cm}ccccc}
\toprule
Region & $W_{\rm fil}$ (pc) & $p$ & $N$ & NN (pc) & $\lambda/W$ \\
\midrule
Taurus  & $0.110$ & $2.01$ & 295 & $0.124$ & $1.1$ \\
Orion B & $0.148$ & $2.00$ & 356 & $0.207$ & $1.4$ \\
Perseus & $0.147$ & $2.00$ & 235 & $0.263$ & $1.8$ \\
Aquila  & $0.161$ & $2.00$ & 446 & $0.210$ & $1.3$ \\
\midrule
median  & $0.148$ & $2.00$ & --- & --- & $1.4$ \\
\bottomrule
\end{tabular}
\par\vspace{2pt}
\footnotesize The measured widths ($0.11$--$0.16$\,pc) agree with the published HGBS values and exceed the canonical $0.10$\,pc, so region-specific Plummer widths give a \emph{lower} $\lambda/W$ ($\approx1.4$, raw) than the raw fixed-width value ($\approx2.1$; Table~\ref{tab:headline}): all four regions are sub-classical under either convention. NN spacings are the full-coverage medians (Section~\ref{sec:injection_recovery}), measured on the same DisPerSE skeleton as the widths (its deposited crest for the width fit, its bridged form for the spacing); the injection--recovery bias correction lowers each $\lambda/W$ by a further $\sim$10--20\%. (Pipeline outputs retained to 2 s.f.; the convention systematic dominates.)
\end{table}

The fit is consistent across the sample: every region returns a Plummer index $p\simeq2.0$, the value found for HGBS filaments generally \citep{Arzoumanian2019}, and a width within $\sim50\%$ of $0.10$\,pc (inter-quartile ranges $\sim0.07$--$0.20$\,pc), so the sub-classical $\lambda/W$ is a property of all four robust regions, not of one anomalous cloud. The single earlier outlier, a spuriously narrow $0.031$\,pc Gaussian crest for Taurus, was a profile-selection effect of fitting the inner ridge rather than the broad Plummer wings (not a pixel-scale artefact: the $3''$ pixels subtend $0.002$\,pc at $135$\,pc); it does not survive the Plummer fit, which gives $0.11$\,pc and $\lambda/W\approx1.1$.

\subsection{Sample Classification and Distance Robustness}
\label{sec:distance_uncertainties}

We classify the eight regions (Table~\ref{tab:sample}) as \textbf{robust} (Orion~B, Aquila, Perseus, Taurus) or \textbf{limited} by three criteria: $N>500$ cores; a well-established distance; and a skeleton simple enough for unambiguous core-to-filament association. Limited regions fail at least one (Ophiuchus on skeleton complexity, the others on small samples) and are excluded from the primary NN measurement. The largest distance revisions (Serpens $+76\%$, Aquila $+68\%$) are independently corroborated: VLBI maser parallaxes validate Gaia DR3 to $<5\%$ where available \citep{Kounkel2018,Xu2021,OrtizLeon2018,OrtizLeon2017}, and Serpens' 458\,pc is supported by extinction estimates \citep[$440$--$460$\,pc;][]{Yan2022,Green2024} and co-spatiality with Aquila \citep{Zucker2020}. Serpens is \emph{limited} and excluded, so any residual blending would only reduce an effective sample size we already take as four. The sub-classical conclusion is independent of these regions: leave-one-out shifts the weighted mean by $<6\%$, reverting the limited regions to their original distances gives a pairwise-median $\lambda/W=2.68$, and dropping Aquila leaves the region-Plummer median at $1.40$ (from $1.35$) and the fixed-width characteristic at $2.1$, both sub-classical.


\subsection{Choice of spacing statistic: the $L/3$ property of the pairwise median}
\label{sec:statistics}

We adopt the nearest-neighbour (NN) spacing as the primary measure of the fragmentation wavelength, reporting the all-pairs median only for comparability with prior HGBS work. The two differ fundamentally because the pairwise median has an analytic $L/3$ convergence property. For $N$ points uniformly distributed on a filament of length $L$, a single pairwise distance $D=|X_i-X_j|$ has the triangular PDF $f(d)=2(L-d)/L^2$, with mean $\langle D\rangle=L/3$ and median $d^\star=L(1-1/\sqrt{2})\approx0.293\,L\approx0.88\,(L/3)$. The median of all $N(N-1)/2$ pairwise distances converges to $d^\star$ as $N\to\infty$, for a random or a perfectly periodic distribution alike, so it measures the extent, not the fragmentation wavelength. In practice it tracks $\sim L/3$ for any filament with $N\gtrsim5$ cores (Monte-Carlo over $N=5$--$500$ gives median$/L\approx0.293\,(1+0.43/N)$, $\le2\%$ inflation for $N\gtrsim20$) and never approaches the fragmentation wavelength ($\sim30$--$100\times$ smaller); for branched skeletons it measures $\approx0.3\times$ the component's arclength extent, and for a realistic length distribution it tracks the rms-length$/3$ (full tables in the Zenodo archive). The NN statistic is unaffected by this bias and is beam-safe: the $18.2''$ HGBS beam corresponds to $0.012$--$0.039$\,pc, $\gtrsim6\times$ smaller than the measured NN spacing.

We reconstruct the along-skeleton core ordering for the four robust regions from the deposited DisPerSE skeleton maps, bridging fragmentation by morphological closing, and extend the NN measurement to all four at full coverage in Section~\ref{sec:injection_recovery}; the limited regions are excluded from the NN measurement owing to branching skeletons or small samples. Figure~\ref{fig:spacing} shows the per-region spacings with revised distances.

\begin{figure*}
\includegraphics[width=0.95\columnwidth]{figures/figure1_spacing_comparison.pdf}
\caption{Nearest-neighbour core spacing $\lambda/W_{\rm fil}$ per robust region, in the two DisPerSE width conventions: directly-measured region-Plummer widths (blue circles, preferred; per-region raw $1.12,1.40,1.79,1.31$, median $1.4$, dashed) and the fixed $0.10$\,pc width (orange squares; $1.2,2.1,2.6,2.1$), both shown \emph{raw} (before the injection--recovery correction; Section~\ref{sec:injection_recovery}). The red band is the classical fastest-growing wavelength $\lambda_{\rm m}\approx5$--$6\times$FWHM (\citealt{Fischera2012}; $=4\times$ the flat diameter, \citealt{Inutsuka1992}) and the dotted line the classical \emph{critical} wavelength $\lambda_{\rm cr}\approx2.6\times$FWHM; every region lies well below both. The grey band marks the independent FilFinder skeleton range ($\lambda/W\approx0.4$--$0.9$).}
\label{fig:spacing}
\end{figure*}

\subsection{Injection--recovery validation and model dependence of the correction}
\label{sec:injection_recovery}

To test whether the NN statistic recovers known spacings under realistic HGBS conditions, we placed synthetic cores with known periodic spacings (0.15--0.35\,pc) on HGBS-property-matched DisPerSE skeletons for all 8 regions. Across 40 tests the raw measurement recovers the input with a small intrinsic-estimator bias of $3.05\pm0.69$\% on clean skeletons; the larger correction actually applied is the \emph{full-coverage pipeline} bias ($\div1.11$--$1.29$, from junction/branch topology on the real bridged skeletons), which moves raw $2.1\to1.9$. Because that test injects a periodic pattern whereas the real cores are clustered (Section~\ref{sec:clustering}), we distinguish two adjustments. The \emph{topological} pipeline bias ($\div1.11$--$1.29$) is a genuine, region-measured correction and is applied. The \emph{point-process model factor} (periodic $\div1.20$ vs clustered $\div1.64$) is not a bias on a line-density observable: if the cores are clustered the measured NN median already \emph{is} the physical quantity, so correcting it toward a periodic template the data reject would be circular. We therefore carry the \emph{raw} value as the adopted $(\lambda/W)_0$ (Section~\ref{sec:injection_recovery_model}) and quote the clustered branch only as a lower bound: applying $\div1.64$ maps raw $\lambda/W\approx2.0$ to $\approx1.2$ (range $0.8$--$1.7$), the periodic factor to $\approx1.9$. The sub-classical \emph{sign} holds under every choice.

Repeating the injection--recovery for the pairwise-median statistic confirms that it does not recover $\lambda_{\rm true}$: it returns a value $\sim$$20\times$ larger, tracking $L/3$ ($\langle\lambda_{\rm PW}\rangle/\lambda_{\rm true}\approx L/(3\lambda_{\rm true})$) and asymptoting to $\approx0.88\,(L/3)$, while the NN statistic recovers $\lambda_{\rm true}$ to $1.8$--$3\%$ on the identical skeletons. Computing the pairwise median over \emph{all} cores in each robust region (without per-filament segmentation) returns $\sim L_{\rm region}/3$ ($12.5$, $7.5$, $8.0$, $2.5$\,pc for Orion~B, Aquila, Perseus, Taurus), the region extent, not the $\sim0.2$--$0.35$\,pc per-filament values, confirming directly that the statistic measures segmentation scale rather than the fragmentation wavelength.

\textbf{Full-coverage along-skeleton NN.} We extend the NN measurement to all four robust regions using the dense deposited DisPerSE skeletons, which (unlike the sparse FilFinder spines) pass near essentially all cores. We bridge the fragmented skeletons by morphological closing plus 1-pixel re-skeletonization, order cores along each filament by graph arclength, and associate cores within a half-width ($\sim$0.06\,pc) of a spine. This yields $N_{\rm pairs}=461,614,563,168$ spacings (Taurus, Orion~B, Perseus, Aquila) with median NN spacings $0.124,0.207,0.263,0.210$\,pc, i.e.\ $\lambda/W=1.2,2.1,2.6,2.1$ (characteristic $\approx2.1$), robust to the bridging and association radii ($\lambda/W$ for Orion~B stays $1.9$--$2.7$). Injecting synthetic cores onto the real bridged skeletons and recovering them shows a $\sim$$15$--$20$\% overestimate (recovery ratio $1.11$--$1.29$, from longest-path ordering at junctions). The periodic pipeline factor gives $\lambda/W\approx1.0,1.9,2.1,1.8$ (characteristic $\approx1.9$), a clustered factor ($\div1.64$) $\approx1.2$; applied uniformly to every width convention (Section~\ref{sec:widths}) the fixed-width characteristic is $\lambda/W\approx1.2$--$1.9$ and the region-Plummer $\approx0.8$--$1.4$. We adopt the raw region-Plummer $(\lambda/W)_0\approx1.4$ and carry the correction as a uniform systematic; its consequence is that the sign of the observation--simulation offset itself depends on the uncertain correction (single-filament fragmentation suffices at $\lambda/W\approx1.2$--$1.4$ but is insufficient at the clustered $\approx0.8$), reinforcing the under-determination conclusion. \textbf{The full-coverage NN spacing ($\lambda/W\approx1.2$--$2.1$) is sub-classical by a factor $\approx3$--$4$ (DisPerSE; more under FilFinder), while the pairwise median $\lambda/W\approx2.8$ overestimates it by $\sim$45--50\% through the $L/3$ bias.} Taurus is the most sub-classical region, not an outlier.

As an independent cross-check, running the same along-skeleton BFS-arclength NN pipeline on independently extracted FilFinder spine maps (all four robust regions) gives $\lambda/W\approx0.4/0.6/0.7/0.8$ for Taurus, Orion~B, Perseus, Aquila respectively, systematically $2$--$3\times$ smaller than the DisPerSE values ($1.0/1.9/2.1/1.8$) because the FilFinder mask threshold produces a denser, more branched network. Both methods agree the characteristic spacing is far sub-classical ($\lambda/W\ll4$) in every region; FilFinder if anything strengthens the conclusion.

\smallskip\noindent\textbf{Model dependence of the correction.}\label{sec:injection_recovery_model}
As a direct test of this model-dependence, for each region we derived a correction matched to its own measured clustering (CoV$=1.36,0.83,0.96,0.65$ for Taurus, Orion~B, Perseus, Aquila) by injecting a gamma renewal process at that CoV and recovering it through a deblending pipeline anchored to the periodic ($\div1.20$) and Poisson ($\div1.64$) endpoints. The region-matched factors are $\div1.86,1.53,1.61,1.40$, giving corrected $\lambda/W\approx0.7,1.4,1.6,1.5$ (fixed $W=0.10$\,pc). The pooled factor slightly over-corrects the regular regions and under-corrects Taurus, but the effect on the characteristic is small (mean $\lambda/W$ $1.2\to1.3$): the sub-classical result holds under each region's own matched correction, and $\lambda/W\approx1.9$ (which requires the rejected periodic model) remains excluded. We retain the pooled value while noting it averages a CoV spanning $0.65$--$1.36$.

\subsection{Regularity of the core spacing: periodic, random, or clustered?}
\label{sec:clustering}
Classical fragmentation predicts a \emph{quasi-periodic} chain of cores; fibre projection or a turbulent/accretion-driven cascade predicts a broad, clustered distribution with no single comb. Ordering associated cores by arclength along each spine and measuring adjacent-core gaps, we find a coefficient of variation $\mathrm{CoV}=\sigma/\mathrm{mean}$ of order unity (Taurus $1.36$, Orion~B $0.83$, Perseus $0.96$, Aquila $0.65$; pooled $1.16$; Fig.~\ref{fig:clustering}).

We calibrated the goodness-of-fit tests by parametric bootstrap, resampling from the fitted distribution to build the exact null distribution of the Kolmogorov--Smirnov statistic and so removing the estimated-parameter bias of the standard test. The result is two-sided and unambiguous: the regular ``periodic$+$jitter'' model is rejected overwhelmingly ($p_{\rm boot}\approx0$ in every region), \emph{and} the pure-exponential (Poisson) model is also rejected ($p_{\rm boot}\leq0.002$ in all four regions). The cores are therefore neither periodic nor pure-Poisson. The exponential fails because of a deficit of the smallest separations (1-D Clark--Evans ratio $R>1$), a finite-core/deblending exclusion at the $\sim$beam scale; a shifted-exponential (minimum-separation) model fits best in every region (by AIC), with an inferred hard core of $0.011$--$0.044$\,pc, of order the beam/core diameter.

The robust statistical statement is thus two-sided --- the data reject \emph{both} simple quasi-periodic fragmentation and pure-Poisson randomness. They are consistent with a point process bearing a beam-scale hard-core exclusion (which drives the Poisson rejection and is partly instrumental) plus region-dependent tail clustering (super-Poisson in Taurus, sub-Poisson in Aquila and Orion~B, so ``clustered'' is a population-level, not a universal, description). Either way there is no single characteristic comb, the observational counterpart of the \citet{Clarke2020} result that fragmentation through sub-filaments leaves no single characteristic length.

Two consequences follow. An injection--recovery correction, if applied, should use the clustered rather than the periodic branch (Section~\ref{sec:injection_recovery}); we nonetheless carry the raw value as adopted. And because for a near-exponential gap law the median is $\ln2\approx0.69$ of the mean --- and it is the mean that equals the inverse line density --- the reported median NN under-states the line-density scale by $\sim$30\%, a stated systematic that only deepens the sub-classical sign. The observed spacing is set by the core line-density plus a finite-core exclusion, not by a resonant wavelength: the classical $\lambda_{\rm m}\approx5$--$6\times$FWHM would place $\sim3$--$4\times$ fewer cores along the spine.
\begin{figure}
\centering
\includegraphics[width=0.99\columnwidth]{figures/fig_clustering.pdf}
\caption{Adjacent-core gap distributions along the skeletons of the four robust regions (histograms), with exponential/Poisson (solid) and regular/periodic (dashed) reference curves; each panel gives the coefficient of variation CoV, the 1-D Clark--Evans ratio $R$, and the exponential goodness-of-fit $p$-value \emph{calibrated by parametric bootstrap} ($p_{\rm boot}$; the null distribution of the KS statistic is built by resampling from the fitted exponential, removing the estimated-parameter bias of the standard test). The result is two-sided: the regular/periodic model is rejected overwhelmingly ($p_{\rm boot}\approx0$) in every region, \emph{and} the pure-exponential (Poisson) model is also rejected ($p_{\rm boot}\leq0.002$ in all four), so the cores are neither periodic nor pure-Poisson --- consistent with a clustered process bearing a beam-scale hard-core exclusion (the deficit of small gaps, $R>1$).}
\label{fig:clustering}
\end{figure}

This does not make the comparison with a fragmentation wavelength meaningless, but it fixes its logical status. A fragmentation process still sets a characteristic \emph{mean} core separation --- the number of cores formed per unit length --- even when the resulting point pattern is not a sharp periodic comb; clustering, a finite-core (beam-scale) exclusion, and projection across blended structure then broaden that mean into the observed near-exponential gap distribution. The observed median separation and any single simulated fragmentation wavelength therefore need not, and should not, coincide better than order-of-magnitude: what is physical and convention-robust is that the core \emph{line density} is sub-classical --- roughly three to four times more cores per unit length than the one-core-per-$\lambda_{\rm m}$ of the classical equilibrium prediction. We accordingly read all wavelength-vs-spacing comparisons below as order-of-magnitude consistency of a characteristic scale, not as a match of a resonant wavelength.

\subsection{Combined Systematic Error Budget}
\label{sec:error_budget}

We distinguish three kinds of uncertainty and present them separately: (i) \emph{statistical} (bootstrap/jackknife); (ii) \emph{systematic} (measured offsets, e.g.\ the pipeline bias and fixed-width convention); and (iii) \emph{sensitivity ranges} (bounds on effects we can sign but not measure as Gaussian errors, e.g.\ protostellar migration). Quoted $\pm$ values combine these in quadrature but are \emph{not} formal confidence intervals. The adopted region-Plummer NN result $(\lambda/W)_0\approx1.4$ (raw fixed-width analogue $\approx2.1$) carries four observational systematics: residual pipeline bias ($\sim$10\%; Section~\ref{sec:injection_recovery}, Table~\ref{tab:headline}); the fixed-width convention ($\sim$25\%; the distance-dependent width variation of \citealt{Panopoulou2022}, counted once); protostellar migration ($\sim$7.5\%); and between-region bootstrap scatter ($\sim$4\%). Manual spur removal is a further $\lesssim5\%$ systematic (proxy; Section~\ref{sec:methodology}). The skeleton algorithm is the single largest observational systematic (FilFinder $2$--$3\times$ smaller than DisPerSE), treated separately below. Because the dominant term (width convention) is a \emph{deterministic} offset, we regard the aggregate as a \emph{min--max envelope}, $\lambda/W\approx0.8$--$1.9$ (Table~\ref{tab:headline}), rather than a Gaussian $\sigma$; for economy we also quote $\sigma_{\rm sys}\approx28\%$ (equivalently $\lambda/W=1.9\pm0.5$ fixed-width), with the skeleton factor of two carried separately.

The legacy pairwise-median value $\lambda/W=2.84$ is an observational $L/3$ scale measure (Section~\ref{sec:statistics}); its observational uncertainty is the bootstrap scatter $\pm4.2\%$ ($\lambda/W=2.84\pm0.12$), and we do not quote a $\sigma$-level tension with the classical prediction because the statistic does not measure the fragmentation wavelength. (In pc, $\lambda/W=2.84$ is $0.284$\,pc; the full- and robust-sample core-weighted pairwise means are $0.279$ and $0.290$\,pc respectively, Table~\ref{tab:sample}, differing from $0.284$ only through the weighting scheme.)

\emph{Skeleton-algorithm dependence.} The factor-of-two spread between DisPerSE and FilFinder is the single largest uncertainty, and the two algorithms do not measure the same object: DisPerSE traces persistence-robust density crests (and built the published HGBS catalogues), whereas FilFinder skeletonises an intensity mask into a denser, more branched network that subdivides a filament and shortens the apparent spacing. Their values ($\lambda/W\approx1.9$ and $\approx0.4$--$0.8$) bracket a genuine ambiguity in what constitutes ``the filament'', not a right-and-wrong choice. Both lie far below the classical $\lambda_{\rm m}\approx5$--$6\times$FWHM, so the factor of two moves the \emph{magnitude} but not the \emph{sign}: the sub-classical result is robust to the skeletoniser even though its precise value is not. A matched-sensitivity, common-footing FilFinder-versus-DisPerSE extraction (matched spatial filtering and an agreed filament-segment definition) is the important future work that would quantify this term; we defer it and carry the factor of two as the dominant open observational systematic, larger than the formal statistical error on any comparison with theory. The skeletoniser also selects the interpretation: in matched Plummer units the simulated single-filament wavelength ($\approx1.0$--$1.3$; Section~\ref{sec:validation_endtoend}) matches the region-Plummer values ($1.2$--$1.4$), sits below fixed-width DisPerSE ($1.9$) and above FilFinder ($0.4$--$0.8$), so a denser FilFinder extraction would require the extra compression of hierarchical-fibre projection (Section~\ref{sec:fibre_projection}) while DisPerSE would not --- only a fibre-resolved measurement (Section~\ref{sec:limitations}) can decide.

Table~\ref{tab:errorbudget} collects every significant uncertainty on the comparison. The dominant terms are \emph{methodological}, not statistical: the skeleton algorithm and the width convention each move $\lambda/W$ by up to a factor of two, and it is these, rather than the formal bootstrap error, that limit how precisely the observed and simulated spacings can be compared.

\begin{table}
\caption{Summary of the uncertainties on the $\lambda/W$ comparison. ``Sim.'' entries apply to the simulated wavelength, ``obs.'' to the observed spacing, and ``comparison'' to terms that enter only when the two are matched.}
\label{tab:errorbudget}
\footnotesize
\setlength{\tabcolsep}{3pt}
\begin{tabular}{p{3.6cm}p{1.9cm}p{1.7cm}}
\toprule
Source & Type & Magnitude \\
\midrule
Bootstrap / cloud-to-cloud (obs.) & sensitivity/descriptive & $\sigma\approx0.4$ ($\sim$20\%) \\
NN pipeline bias (obs.) & systematic & $\sim$10\% (corrected) \\
Protostellar migration (obs.) & sensitivity & $\sim$7.5\% \\
Width convention, fixed vs region (obs.) & systematic & $\sim$25\% \\
Skeleton algorithm, DisPerSE/FilFinder (obs.) & systematic & factor $\sim$2 \\
Manual skeleton editing (obs.) & systematic (proxy) & $\lesssim$5\% \\
Growth-rate resolution/epoch (sim.) & modelling & $\lesssim$5\% (converged) \\
Turbulence sensitivity (sim.) & sensitivity & $<$5\% ($\lambda$; single realisation) \\
Line-mass ($f$) dependence (sim.) & modelling & $\sim$25\% \\
Gaussian vs Plummer FWHM (comparison) & systematic & factor $\sim$1.5 \\
\bottomrule
\end{tabular}
\end{table}

When simulation $\lambda/W_{\rm core}$ values are converted to observable $\lambda/W_{\rm fil}$, the T1 width-normalisation (forward-modelled $\eta_W=0.59$--$0.78$, provisional calibration $0.61\pm0.10$; Section~\ref{sec:width_correction}) is the dominant additional systematic; the directly measured supercritical ensemble (Section~\ref{sec:modecompetition}) carries its own $\pm0.4$ scatter. These comparison-only systematics do not affect the NN observational measurement, which is already in $W_{\rm fil}$ units.

Because the four robust clouds are the genuine independent units, we quote uncertainties descriptively, not inferentially. Two distinct quantities must be kept apart: (i) the empirical \emph{per-region range} of the bias-corrected fixed-width values, $\lambda/W=1.0$--$2.1$ (Taurus, Orion~B, Perseus, Aquila $=1.0,1.9,2.1,1.8$; Section~\ref{sec:injection_recovery}), which is the real cloud-to-cloud spread; and (ii) the \emph{methodological/systematic envelope}, $\lambda/W\approx1.1$--$2.4$, obtained by varying the width-convention, skeleton and association-radius choices (the quadrature sum of these is $\sigma_{\rm sys}\approx28\%$, a synthesis of heterogeneous sensitivity ranges, not a Gaussian error). With $N=4$ we do \emph{not} report a bootstrap $\sigma$ or a confidence interval as a formal quantity, and we do not claim a $\sigma$-level exclusion of the classical prediction; the informal descriptor ``$\sigma\approx0.4$'' quoted elsewhere for the four-cloud scatter should be read only as the half-width of the per-region range.

\subsection{NN pipeline reproducibility, population splits and hierarchical statistics}
\label{sec:nnpipeline}

The along-skeleton NN pipeline thresholds the deposited DisPerSE skeleton, bridges fragmentation by morphological closing and re-skeletonisation, associates each core to the nearest spine within $0.06$\,pc, orders cores along each component by graph arclength (double-BFS diameter), and takes adjacent-core separations. It reproduces the deposited catalogue counts and raw medians exactly (Table~\ref{tab:nnpipeline}); cores not on a robust skeleton segment are rejected.

\begin{table}
\caption{NN pipeline counts: skeleton components, catalogue cores, cores associated to a skeleton, and adjacent NN pairs, per robust region.}
\label{tab:nnpipeline}
\footnotesize
\setlength{\tabcolsep}{4pt}
\begin{tabular}{lcccc}
\toprule
Region & comps & catalogue & associated & NN pairs \\
\midrule
Orion B & 67 & 1844 & 680 & 614 \\
Aquila  & 14 & 749  & 182 & 168 \\
Perseus & 54 & 816  & 619 & 563 \\
Taurus  & 32 & 536  & 479 & 461 \\
\bottomrule
\end{tabular}
\end{table}

The fraction of catalogue cores associating to a robust skeleton segment varies across regions --- Orion~B $37\%$, Aquila $24\%$, Perseus $76\%$, Taurus $89\%$ --- so in Orion~B and Aquila the NN spacing is measured on a minority of cores. This reflects skeleton \emph{coverage}: the distributed, clustered clouds have many cores off any robust crest, whereas filament-dominated clouds place almost all cores on a spine. Any bias runs \emph{toward} the sub-classical result: rejected cores lie on fainter, lower-line-density segments with \emph{longer} spacings, so omitting them biases the median \emph{short}. The decisive check is that the result does not depend on the low-coverage regions: restricting to the high-coverage clouds (Perseus, Taurus) gives $\lambda/W=1.8$ and $1.1$, median $1.4$ --- identical to the four-region median. The Orion~B/Aquila medians are also stable to $\lesssim5\%$ as the association half-width varies over $0.04$--$0.08$\,pc, and the prestellar subset matches the all-cores measurement, so the effect is not load-bearing.

To test whether mixing core populations biases the result, we repeated the measurement separately for prestellar, unbound starless, and protostellar cores using the published classification flags. The prestellar-only robust-median $\lambda/W=2.10$ ($717$ pairs) is statistically indistinguishable from the all-cores value $2.08$ ($1806$ pairs; both $\approx1.9$ after the injection-recovery bias correction). Prestellar cores dominate the on-skeleton budget in every region and set the result; unbound starless cores give a larger, noisier value (median $\lambda/W\approx3.3$) and protostellar cores are too few ($6$--$80$ pairs) to constrain it. The sub-classical result is therefore not an artefact of population mixing.

\smallskip\noindent\textbf{Hierarchical statistics and angular-resolution robustness.}

Cores within a filament are not independent, so we assessed the cloud-to-filament-to-core nesting with a three-level nonparametric bootstrap (clouds $\to$ filament components $\to$ core spacings, 2000 iterations, carrying spacings that share a central core together). On a 2-D projected NN metric (which underestimates the arc length and sits below the adopted value, $\lambda/W\approx1.1$--$1.5$) the bootstrap median is stable; the dominant uncertainty is the intrinsic cloud-to-cloud scatter, $\sigma\approx0.4$ in $\lambda/W$ (range $\sim0.8$ in Taurus to $\sim1.7$ in Aquila). Core-count and equal-cloud weighting agree to within $\sim0.1$. We therefore quote the cloud-to-cloud scatter, not a per-pair formal error, as the relevant uncertainty (Section~\ref{sec:error_budget}).

Angular-resolution blending: the HGBS beam grows from $0.012$\,pc (Taurus, 135\,pc) to $0.039$\,pc (Aquila, 436\,pc), and the fraction of nearest-neighbour pairs closer than $2\times$ the physical beam rises with distance ($0$\,per cent in Taurus to $10$\,per cent in Orion~B; Spearman $\rho=+0.80$), as expected. Merging cores within one beam shifts the characteristic $\lambda/W$ by $\le0.004$, however, so distance-dependent blending is not a significant bias on the adopted value.

\begin{table}
\caption{Summary of validation tests performed in this paper.}
\label{tab:validation_summary}
\footnotesize
\setlength{\tabcolsep}{3pt}
\begin{tabular}{p{2.4cm}p{2.8cm}p{2.0cm}}
\toprule
Test & What it checks & Result \\
\midrule
Injection--recovery (NN) & NN bias on real skeletons & $3\%$; pipeline $15$--$20\%$ \\
Injection--recovery (pairwise) & $L/3$ convergence & tracks $L/3$, not $\lambda$ \\
Resolution convergence & $t_{\rm frag}(128^3)$ vs $256^3$ & ratio $0.915\pm0.012$ \\
Truelove criterion & cells per $\lambda_{\rm J}$ at measurement & long.\ $5$--$7$; perp.\ $1$--$2$ \\
IC sensitivity & Gaussian vs uniform profile & identical \\
EOS sensitivity & $\gamma=0.8$--$5/3$ & $\gamma<1$ accelerates; $5/3$ suppresses \\
Region resampling (descriptive) & spread over 4 clouds & $0.261$--$0.298$\,pc \\
Leave-one-out & single-region dominance & $<6\%$ shift \\
Population splits & prestellar vs all cores & $\lambda/W\!=\!2.10$ vs $2.08$ \\
Hierarchical bootstrap & cloud$\to$filament$\to$core & $\sigma_{\rm cloud}\!\approx\!0.4$ \\
Mock-blending & distance-dependent beam & $\Delta\lambda/W\!\le\!0.004$ \\
Region-specific Plummer widths & beam-deconvolved $W_{\rm fil}$ & $0.11$--$0.16$\,pc; $\lambda/W\!=\!1.1$--$1.8$ \\
\bottomrule
\end{tabular}
\end{table}

\subsection{Comparison with Literature}

Our revised-distance measurements differ from the original HGBS values (Table~\ref{tab:literature}) by the region-dependent distance revisions (Aquila $+68\%$; Orion~B essentially unchanged, its pre-Gaia distance already $\sim400$\,pc; Taurus $-4\%$); the core-count-weighted mean rises only $\sim10\%$ for the robust regions, dominated by Aquila.

\begin{table*}
\caption{Comparison with Published HGBS Spacing Measurements}
\label{tab:literature}
\begin{tabular}{lccccc}
\toprule
Region & This Work, NN (pc) & Original HGBS (pc) & This Work, pairwise (pc) & Literature (pc) & Reference \\
\midrule
\textbf{Robust Regions} &  &  &  &  &  \\
Orion B & \textbf{0.207} & 0.212 & 0.313 $\pm$ 0.047 & 0.20--0.30 & \citet{Konyves2020} \\
Aquila  & \textbf{0.210} & 0.206 & 0.346 $\pm$ 0.047 & 0.22--0.26 & \citet{Konyves2015} \\
Perseus & \textbf{0.263} & 0.199 & 0.248 $\pm$ 0.040 & 0.22 $\pm$ 0.02 & \citet{Pezzuto2021} \\
Taurus  & \textbf{0.124} & 0.206 & 0.198 $\pm$ 0.040 & 0.19 $\pm$ 0.02 & \citet{Hacar2013} \\
\midrule
\textbf{Limited Regions} &  &  &  &  &  \\
Ophiuchus & --- & 0.197 & 0.206 $\pm$ 0.053 & 0.18 $\pm$ 0.02 & \citet{Konyves2020} \\
Serpens  & --- & 0.188 & 0.331 $\pm$ 0.097 & 0.17--0.19 & \citet{Konyves2020} \\
TMC1     & --- & 0.203 & 0.195 $\pm$ 0.056 & $\sim$0.20 & \citet{Hacar2013} \\
CRA      & --- & 0.212 & 0.248 $\pm$ 0.072 & $\sim$0.21 & \citet{Konyves2020} \\
\bottomrule
\end{tabular}

\vspace{0.1in}
\footnotesize
\emph{Notes} (1) The primary statistic is the \emph{nearest-neighbour} (NN) spacing (bold): the homogeneous re-analysis leaves Orion~B and Aquila essentially unchanged, \emph{shortens} Taurus ($0.124$ vs $0.206$\,pc; highest association fraction, so most reliable), and \emph{lengthens} Perseus ($0.263$ vs $0.199$\,pc), so it does \emph{not} inflate the spacing the way the pairwise median does, and the net effect strengthens the sub-classical result. The ``This Work, pairwise'' column is retained only as a legacy comparator; it is inflated towards $L/3$ (Section~\ref{sec:statistics}) and is not a physical spacing. NN spacings for the limited regions are omitted. Perseus is the one region where the NN ($0.263$) slightly exceeds the pairwise median ($0.248$), because arclength ordering on its heavily bridged skeleton inflates path distances; both remain sub-classical. (2) ``Original HGBS'' values use pre-Gaia distances. (3) Literature ranges reflect different methods and distance assumptions. (4) Serpens used $260$\,pc; revised $458$\,pc ($+76\%$). (5) For Orion~B the pre-Gaia distance was already $\sim$400\,pc \citep{Konyves2020}, so the pairwise $0.212\to0.313$\,pc difference is a statistic-definition effect, not a distance revision.
\end{table*}

We decomposed this shift on the real Orion~B data by toggling each methodological choice at fixed distance. The dominant term by far is the \emph{statistic definition}: switching from NN to the pairwise median inflates the spacing by $\approx1.5$ ($0.207\to0.313$\,pc; the pairwise median is biased towards $L/3$, Section~\ref{sec:statistics}), accounting for essentially the entire $0.212\to0.313$\,pc difference in Table~\ref{tab:literature}. Skeleton re-extraction ($+4\%$), association half-width ($\pm12\%$) and 2-D$\to$arclength ordering ($+12\%$) each contribute at the $\le15\%$ level and largely cancel; the distance revision ($400\to386$\,pc) contributes $-3.5\%$. On our \emph{primary} NN statistic Orion~B is $0.207$\,pc, within $3\%$ of \citet{Konyves2020}, so the apparent shift is a property of the pairwise-median comparability statistic, not the physical spacing. The like-for-like NN re-analysis is not uniform: it leaves Orion~B and Aquila essentially unchanged, \emph{shortens} Taurus ($0.124$ vs $0.206$\,pc; highest association fraction, so most reliable), and \emph{lengthens} Perseus ($0.263$ vs $0.199$\,pc), so the net effect if anything strengthens the sub-classical conclusion.

\section{THEORETICAL FRAMEWORK}
\label{sec:theory}

\subsection{Classical Fragmentation Theory}

For an isothermal, self-gravitating cylinder of density $\rho_0$ and sound speed $c_s$, the critical line mass for gravitational instability is \citep{Ostriker1964}:
\begin{equation}
\mu_{\rm crit} = \frac{2c_s^2}{G}
\end{equation}
The most unstable fragmentation wavelength is, as an \emph{exact} result for the isothermal equilibrium cylinder \citep{Inutsuka1992},
\begin{equation}
\lambda_{\rm IM92} = 22\,H = 4\times(2R_{\rm flat}),
\label{eq:im92}
\end{equation}
where $H = c_s\,(4\pi G\rho_c)^{-1/2}$ is the thermal scale height, $R_{\rm flat}=\sqrt8\,H$ the flattening radius, and $W_{\rm fil}\approx0.1$\,pc the characteristic filament FWHM. Solving the dispersion relation exactly (Appendix~\ref{sec:val_appendix}; \citealt{Fischera2012}) gives a fastest-growing wavelength $\lambda_{\rm m}\approx5\times W_{\rm fil}$ near-critical, rising to $\approx6\times W_{\rm fil}$ at low line mass, and a critical wavelength $\lambda_{\rm cr}\approx2.6\times W_{\rm fil}$. The often-quoted ``$\lambda\approx4\times$ the width'' is $4\times$ the flat \emph{diameter} $2R_{\rm flat}$ \citep{Inutsuka1992}, which equals $\approx5$--$6\times$ the \emph{FWHM}; taking it as $4\times$ the FWHM understates the classical scale. We therefore adopt $\lambda_{\rm m}\approx5$--$6\,W_{\rm fil}$ (Fischera \& Martin 2012) as the benchmark. This ratio is defined on the equilibrium ($p=4$) FWHM whereas observed filaments have $p\simeq2$; because we compare in FWHM units and forward-model the simulations to a beam-convolved Plummer FWHM (Section~\ref{sec:width_correction}), the sub-classical deficit does not depend on the profile-index difference. We compare as the density-independent ratio $\lambda/W_{\rm fil}$, sidestepping that $22H\propto\rho_c^{-1/2}$ is not fixed in absolute units. The deficit is not new in sign: \citet{Fischera2012} noted observed cores sit ``about the FWHM'' apart, and \citet{Zhang2023cal} find a spacing $\approx1.3\times$ the width in California; our contribution is to quantify it homogeneously with Gaia DR3 and reproduce it in a controlled MHD campaign.

With $f=\mu_{\rm line}/\mu_{\rm crit}$ and $t_{\rm frag}\propto(f-1)^{-1/2}$, marginally supercritical filaments fragment slowly. At magnetically supercritical field strength ($\beta\gtrsim0.2$) all configurations fragment for $\mathcal{M}\geq1$ at every $f$ tested once integrated adequately: above the critical line mass there is no \emph{thermal} stability boundary. This is distinct from \emph{magnetic} stabilisation, where a strong longitudinal field ($\beta\lesssim0.15$) suppresses collapse even for $f>1$ (the ``strong-field, magnetically supported'' regime; Appendix~\ref{sec:supplementary_campaigns_app}); the thermal ($f$) and magnetic (mass-to-flux) criticalities are physically distinct.

We adopt a single reference $W_{\rm fil}=0.10$\,pc for all regions, for comparability with prior HGBS work. The universality of this width is debated \citep{Panopoulou2017,Panopoulou2022}. We have measured region-specific Plummer widths directly (Section~\ref{sec:widths}); they are $0.11$--$0.16$\,pc ($p\simeq2.0$) and give $\lambda/W\approx1.4$, so the sub-classical conclusion is robust to the width convention, and the width term is carried once as the single $\sim25\%$ systematic in the Section~\ref{sec:error_budget} budget.

\subsection{Width normalisation: $W_{\rm core}$ to $W_{\rm fil}$}
\label{sec:width_correction}
Three widths enter the simulation--observation comparison and must be kept separate (Table~\ref{tab:widthdefs}): the initial Gaussian half-width $W_{\rm core}=0.3\,\lambda_{\rm J}$; the formation FWHM $W_{\rm form}=0.683\,\lambda_{\rm J}$ of the evolved projected profile ($\simeq2.28\,W_{\rm core}$, slightly below the initial $2.355\,W_{\rm core}$ owing to contraction; resolution- and seed-converged to $<1\%$); and the observed HGBS width $W_{\rm fil}$, the beam-convolved Plummer FWHM \citep{Arzoumanian2019}, which we obtain by forward-modelling the simulated filaments through the full synthetic-HGBS pipeline (projection, beam convolution, Plummer fit). This gives $\eta_W=W_{\rm form}/W_{\rm fil}=0.59$--$0.78$ and
\begin{equation}
T_1 \equiv \frac{W_{\rm form}}{W_{\rm fil}} = 0.61 \pm 0.10 \quad (\text{forward-model band } 0.59\text{--}0.78),
\label{eq:t1_correction}
\end{equation}
so that $W_{\rm fil}=W_{\rm form}/T_1\approx1.12\,\lambda_{\rm J}$. The observable is $\lambda/W_{\rm fil}$; because the fragmentation spacing $\lambda$ is a physical length independent of the width definition, the conversion from the simulation-internal ratio $\lambda/W_{\rm core}$ (normalised by the initial $\sigma=0.3\,\lambda_{\rm J}$) is
\begin{equation}
\left(\frac{\lambda}{W}\right)_{\rm fil} = \frac{W_{\rm core}}{W_{\rm fil}}\left(\frac{\lambda}{W}\right)_{\rm core} = T_1\,\frac{W_{\rm core}}{W_{\rm form}}\left(\frac{\lambda}{W}\right)_{\rm core} \approx 0.27\left(\frac{\lambda}{W}\right)_{\rm core},
\label{eq:t1_apply}
\end{equation}
i.e.\ $\lambda/W_{\rm fil}=\lambda\,[\lambda_{\rm J}]/1.12$, the spacing in units of the forward-modelled observed width; the factor $0.27$ includes the $\sigma\!\to\!$FWHM factor $0.44$, and using $T_1$ alone would overstate $\lambda/W_{\rm fil}$ by $\approx2.3\times$. The observational NN measurement (Section~\ref{sec:injection_recovery}) is already in $W_{\rm fil}$ units. Because the observable is sensitive to how $W_{\rm fil}$ is defined, we do not rely on it; the primary comparison uses the convention-free growth-rate wavelength (an upper limit; Section~\ref{sec:validation_endtoend}), and this normalisation is retained only as a cross-check.

\begin{table}
\caption{The three filament widths used in the simulation--observation comparison (all in Jeans lengths, $\lambda_{\rm J}=1$).}
\label{tab:widthdefs}
\footnotesize
\setlength{\tabcolsep}{3pt}
\begin{tabular}{p{1.5cm}p{1.4cm}p{4.0cm}}
\toprule
Symbol & Value ($\lambda_{\rm J}$) & Definition \\
\midrule
$W_{\rm core}$ & $0.30$ & initial Gaussian half-width $\sigma$ (input) \\
$W_{\rm form}$ & $0.683$ & Gaussian FWHM of evolved projected profile, supercritical ensemble ($\simeq2.28\,W_{\rm core}$) \\
$W_{\rm FWHM}$ & $0.58$ & Gaussian FWHM at the growth-rate epoch, near-critical runs (narrower, already contracting); normalises Table~\ref{tab:wave} \\
$W_{\rm fil}$ & $\approx1.12$ & beam-convolved Plummer FWHM ($=W_{\rm form}/T_1$); the observable normalisation \\
\bottomrule
\end{tabular}
\end{table}

For the ambient-confined supercritical ensemble ($\lambda/W_{\rm core}=3.27$) this gives $\lambda/W_{\rm fil}\approx0.9$.

\subsection{Magnetic Tension Mechanism}
\label{sec:magnetic}

For a filament threaded by a magnetic field at an angle $\theta$ to the axis (Alfv\'en speed $v_A$), the general oblique dispersion relation is \citep{Nakamura1993}:
\begin{equation}
\omega^2 = c_s^2 k^2 + v_A^2 k^2 \sin^2\theta - 4\pi G\rho_0 F(kR)
\label{eq:dispersion}
\end{equation}
where $F(kR)$ encodes the cylindrical geometry and $\theta$ is the field--axis angle ($0^\circ$ longitudinal, $90^\circ$ perpendicular). This is the thin-perturbation \emph{static-field} relation: it treats $\mathbf{B}$ as a fixed background and does \emph{not} describe a current-carrying (helical) equilibrium, which is why the helical case cannot be read off Eq.~(\ref{eq:dispersion}) and needs the separate treatment of Section~\ref{sec:optional_tests}. For the axial mode of a \emph{longitudinal} field the compression is along $\mathbf{B}$, so the magnetic term $v_A^2k^2\sin^2\theta$ vanishes and the relation reduces to the field-free thermal one: a static longitudinal field leaves the wavelength almost unchanged (the term is maximal for $\theta=90^\circ$). We do \emph{not} solve the full cylindrical dispersion relation, so this paper contains \emph{no quantitative model of magnetic tension}. Consistent with this, our longitudinal runs show the growth-rate wavelength depends only weakly on $\beta$ ($\lambda_{\rm max}=1.00,1.14,0.80\,\lambda_{\rm J}$ at $\beta=0.3,0.5,1.0$, within the mode-quantisation step $\Delta\lambda\approx0.15\,\lambda_{\rm J}$, not a resolved trend) and on supercriticality, so the scale is set by the thermal Jeans length, not the field. The shortening we measure is therefore \emph{not} magnetic tension but the dynamical/contraction shift; a genuinely shorter magnetic wavelength would need a force-balanced \emph{helical}/toroidal field, which remains untested (Section~\ref{sec:optional_tests}). Planck finds dense filaments preferentially perpendicular to the mean \emph{projected} field \citep{Jadhav2026}; the projected orientation does not fix the internal geometry, so the internally-longitudinal fraction is unconstrained (a minority, $\lesssim10\%$, if it tracks the projection), the assumption under which our longitudinal result applies to part of the population.

\subsection{Hierarchical Fragmentation}

Filament fragmentation is hierarchical: filaments fragment into velocity-coherent fibres, and fibres into cores \citep{Hacar2013,Hacar2018,Yang2024}. Whether the fibre-to-core spacing is quasi-periodic at the classical scale is not established --- the ALMA study of \citet{Yang2024} finds a clear filament--fibre--core hierarchy but only a weak quasi-periodic imprint at the core level. If HGBS filaments are bundles of fibres each fragmenting near the classical scale, measuring core spacings across the whole filament compresses the apparent spacing, in the observed direction. We test this quantitatively in Section~\ref{sec:fibre_projection}, where a forward-model of $N$-fibre bundles run through our NN pipeline reproduces the observed $\lambda/W$; it emerges as a viable but not unique explanation (single-filament dynamical fragmentation is also consistent), and confirming it requires fibre-resolved core-spacing analysis beyond the single Orion~B region of \citet{Yang2024}.

\section{MHD SIMULATIONS}

The full campaign comprises $2251$ Athena++ runs across 18 grids (Table~\ref{tab:campaigns}); it is a set of controlled numerical experiments, not models of individual HGBS filaments. Most of these runs establish \emph{outcomes and timescales} across parameter space; the headline fragmentation-\emph{wavelength} measurement rests on the 25-run near-critical growth-rate campaign (whose dominant-mode measurements across $f$, $\beta$ and resolution are reported in Table~\ref{tab:wave}; the 8 isotropic $128^3$--$512^3$ runs at $f=1.1$ form the production convergence ladder), so the large campaign size should not be read as statistical weight behind the wavelength. Two results emerge (Section~\ref{sec:why_mhd_fails}). The \emph{outcome} of a run, radial collapse or longitudinal beading, depends on regime and boundary conditions, so a single filament has no setup-independent fate. Its dominant fragmentation \emph{mode}, measured from the growth-rate spectrum, is sub-classical ($\approx0.5$--$1\times$ the classical value; Section~\ref{sec:validation_endtoend}) for near-critical, longitudinal-field filaments; the perpendicular geometry is not numerically resolved. This section provides the evidence for both.

\subsection{Simulation Methodology}

We performed self-gravitating MHD simulations using Athena++ \citep{Stone2020} to test filament fragmentation under controlled conditions. The simulations solve the ideal MHD equations with self-gravity:
\begin{align}
\frac{\partial \rho}{\partial t} &= -\nabla \cdot (\rho \mathbf{v}) \\
\frac{\partial (\rho \mathbf{v})}{\partial t} &= -\nabla \cdot \left(\rho \mathbf{v} \mathbf{v} + P + \frac{B^2}{8\pi} - \frac{\mathbf{B}\mathbf{B}}{4\pi}\right) - \rho \nabla \phi \\
\frac{\partial \mathbf{B}}{\partial t} &= \nabla \times (\mathbf{v} \times \mathbf{B}) \\
\frac{\partial E}{\partial t} &= -\nabla \cdot \left[(E + P + \frac{B^2}{8\pi})\mathbf{v} - \frac{(\mathbf{v} \cdot \mathbf{B})\mathbf{B}}{4\pi}\right] - \rho \mathbf{v} \cdot \nabla \phi
\end{align}
where $\rho$ is density, $\mathbf{v}$ is velocity, $\mathbf{B}$ is magnetic field, $P$ is pressure, $E$ is total energy, and $\phi$ is gravitational potential satisfying $\nabla^2 \phi = 4\pi G\rho$.

We initialise a cylindrical filament with density profile $\rho(r)=\rho_0/[1+(r/W)^2]$ and uniform field $\mathbf{B}_0$, characterised by the thermal line-mass fraction $f=M_{\rm line}/M_{\rm line,crit}$ (Eq.~2), plasma $\beta=8\pi P/B^2$, and the seed-amplitude parameter $M_{\rm seed}=\delta v/c_s\times10^{4}$. At filament conditions ($n\sim10^4\,{\rm cm^{-3}}$, $T\sim10$\,K) $\beta=1$ corresponds to $B\sim15\,\mu$G and $\beta=0.3$ to $\sim30\,\mu$G, the $\mu$G--tens-of-$\mu$G range of dense-filament estimates \citep{Crutcher2012}. $M_{\rm seed}$ is a perturbation-spectrum label, not a physical Mach number: the base grid uses a single-mode seed $\delta v/c_s=10^{-4}$, whereas the growth-rate ladder of Section~\ref{sec:validation_endtoend} uses a $10^{-3}$ twelve-mode broadband seed. The magnetic mass-to-flux ratio is distinct from the thermal $f$; with $B_0$ set by $\beta$ and the central density by $f$, it scales as $\mu_\Phi/\mu_{\Phi,\rm crit}\propto f\sqrt{\beta}$, and every configuration has $\mu_\Phi/\mu_{\Phi,\rm crit}>1$ (none strictly magnetically subcritical), so the low-$\beta$ runs are magnetically \emph{supported}.

We explored a base grid spanning $\mathcal{M}=0.5$--$5$ and $\beta=0.05$--$5$ at a near-critical line mass, and a transition grid across $f=1.4$--$2.0$, $\beta=0.3$--$1.3$ and $\mathcal{M}=1$--$5$ to resolve the stable--unstable boundary. Full ranges and run counts are in Appendix~\ref{sec:campaigns} (Table~\ref{tab:campaigns}).

In brief (full detail in Appendix~\ref{sec:supplementary_campaigns_app}), the base grid establishes three outcome regimes: magnetically supported quasi-stability at strong fields ($\beta\lesssim0.15$), radial collapse, and longitudinal fragmentation; the near-critical, longitudinal-field case that is our focus beads at $t_{\rm frag}\approx1.1$--$1.2\,t_{\rm J}$, and turbulence changes the fragmentation \emph{wavelength} by $<5\%$ while shortening the timescale $3$--$5\times$. The radial-collapse/fragmentation timescale competition, the three-regime $(\beta,\mathcal{M})$ map, and the turbulence/critical-transition dependence are given in Appendix~\ref{sec:supplementary_campaigns_app}.


\subsection{Ambient confinement and the direct supercritical measurement}
\label{sec:modecompetition}

The outcome for a supercritical filament is set by whether it is confined by an ambient medium, as real HGBS filaments are. Throughout, $d$ denotes the transverse distance from the filament axis to the domain boundary in Jeans lengths. The main 654-run grid uses a truncated Gaussian with no ambient medium, so the Gaussian skirt wraps through the periodic boundary and feeds an accretion channel that suppresses beading, and the filament collapses radially (Appendix~\ref{sec:regimes}). With ambient confinement (\texttt{filament\_ambient} generator), the same configurations fragment longitudinally: a reconciliation suite (12 runs, $f\times\beta\times2$ seeds) beads at physical, $\beta$-dependent collapse times ($0.55$--$0.70\,t_{\rm J}$), converging across boundary conditions.

Ambient-confined supercritical filaments ($f=1.5$--$2.0$, $\theta=0^\circ$) fragment longitudinally: an ensemble of 39 runs (periodic and reflecting BCs, ambient confinement) develops 3--10 interior axial density maxima at a spacing of order the Jeans length. We do not rely on the median peak spacing (counting peaks late in the collapse is epoch- and resolution-dependent; Section~\ref{sec:validation_endtoend}); we take the wavelength from the linear growth-rate spectrum, whose low-$n$ modes are resolution-stable, giving $\lambda_{\rm max}\approx1.0$--$1.14\,\lambda_{\rm J}$, $\approx0.35$--$0.5$ times the classical value, declining mildly with $f$ and $\beta$. The ambient medium confines the filament (central-to-ambient density ratio $\sim10^2$, in approximate radial force balance, consistent with HGBS crest-to-background contrasts of $10$--$10^2$; \citealt{Arzoumanian2019}), so it is representative rather than tuned. The 39 runs span periodic ($n=12$, median $\lambda/W_{\rm fil}\approx0.89$), reflecting ($n=26$, median $\approx1.01$) and one user-outflow boundary; periodic and reflecting are statistically indistinguishable (KS $p=0.17$). This ensemble median ($\approx0.9$--$1.0$) and the growth-rate result are the same physical wavelength, $\lambda\approx1.1\,\lambda_{\rm J}$, in different width normalisations. The ensemble beads early ($t_{\rm bead}\approx0.6$--$0.75\,t_{\rm J}$), before runaway and at moderate contrast, so with $\lambda\approx\lambda_{\rm J}$ resolved by ${\sim}32$ cells it satisfies Truelove at measurement (Appendix~\ref{sec:validation}). The full per-run table is released with the public archive.

Perpendicular-field filaments ($\theta=90^\circ$) do not universally spindle (smooth radial/axial collapse without discrete beading). An 18-run ideal-MHD grid at $d=1.0$ shows $7/18$ bead (at $\beta=0.5$ and $2.0$), with spindle collapse concentrated near equipartition ($\beta\approx1$), reflecting a competition between the axial and radial modes (a plausibility argument, not a mode-by-mode demonstration). Where they bead, the spacing appears short ($\lambda/W_{\rm fil}\approx0.4$) --- but this \emph{conflicts} with linear theory, which predicts a \emph{longer} mode for a static perpendicular field; since these runs are unresolved ($N_J\approx1$--$2$) we draw no perpendicular conclusion here (a transverse-refinement ladder to $N_J\approx170$ resolves it in Section~\ref{sec:optional_tests}: the ideal perpendicular filament spindles and the short beading is a Truelove artefact). In an expanded 36-run non-ideal survey ($\eta_{\rm AD}=0.001$--$0.1$, $\beta=0.5$--$2.0$), $32/36$ bead, but at $N_J\approx1$--$2$ where under-resolution itself produces spurious fragmentation \citep{Truelove1997}, so we cannot separate a genuine ambipolar enhancement from a numerical one. The median onset wavelength is $\lambda/W_{\rm fil}\approx0.6$ (bulk $0.6$--$0.7$; strongest $0.85$). These are fragmentation-onset values: the runs bead at $t\approx0.44$--$0.55\,t_{\rm J}$ and reach runaway ($C\sim10^4$) within $\Delta t\approx0.02$--$0.04\,t_{\rm J}$, and a finer-cadence rerun shows no converged plateau. Ambipolar diffusion therefore promotes but does not resolve the perpendicular tension; a converged measurement is left to future work.

\subsection{Testing the resolution and magnetic-field alternatives}
\label{sec:optional_tests}
Three targeted campaigns test the concerns that most affect the interpretation: whether the longitudinal wavelength is resolution-limited, whether the perpendicular geometry produces genuine short beading, and whether a helical field --- the one configuration expected to shorten the axial mode --- actually does.

\emph{Resolution stability of $\lambda_{\rm max}$.} Two refinement directions behave oppositely. Under \emph{transverse} refinement (the isotropic $128^3$--$512^3$ ladder) the short-wavelength spectrum ($\lambda\lesssim0.9\,\lambda_{\rm J}$, $n\gtrsim9$) grows with resolution and is numerical (unresolved transverse sub-fragmentation; Fig.~\ref{fig:growthrate}a), whereas the physical band ($\lambda\gtrsim1\,\lambda_{\rm J}$, $n\leq8$) is converged. Under \emph{axial-only} refinement ($512\times128^2\to1024\times128^2$) the physical band reproduces to better than $0.5\%$ ($\Gamma$ at $n=6,7,8$, i.e.\ $\lambda=1.33,1.14,1.00\,\lambda_{\rm J}$, is $17.6,18.7,18.7$ at both resolutions, plateauing without turning over), \emph{and} the short-$\lambda$ rise is unchanged, confirming it is fixed by the axial Nyquist scale, not a physical mode migrating to shorter $\lambda$. The $\lambda_{\rm max}\lesssim1\,\lambda_{\rm J}$ upper bound is therefore stable: axial refinement does not move the physical-band plateau, and the only shorter-$\lambda$ structure is the transverse-resolution artefact already identified.

\emph{The perpendicular geometry: a resolved spindle, not short beading.} On a transverse-refinement ladder ($128^3$--$512^3$) the near-critical \emph{ideal} perpendicular filament collapses monolithically in the radial direction at every resolution, the axial mode staying at the box scale and the axial perturbation at seed level throughout the well-resolved phase (at $512^3$ the spine stays resolved with $N_J\approx170$ cells at $C\approx3$, falling only to $N_J\approx40$ at $C\approx30$ --- still $\approx10\times$ above Truelove; Appendix~\ref{sec:validation}). Short-wavelength axial structure ($\lambda/W\approx0.4$) appears only once the spine drops below $N_J\approx1$; refinement delays its onset to deeper collapse (Fig.~\ref{fig:perp_ladder}), the signature of a grid artefact. The perpendicular field therefore \emph{suppresses} axial fragmentation in favour of radial collapse, as static linear theory predicts (a perpendicular field lengthens the axial mode), and the short low-resolution beading is numerical. This is for \emph{ideal} MHD; the ambipolar-perpendicular case (Section~\ref{sec:modecompetition}) remains separately unconverged.

\emph{A seeded helical field pinches the width but does not shorten $\lambda_{\rm max}$; the force-balanced case remains untested.} We added a divergence-free helical field (uniform axial $B_z$ plus a toroidal $B_\phi(r)$ peaking at $r_c=W_{\rm core}$, from a discrete curl of $A_z(r)=B_t r_c\ln[1+(r/r_c)^2]$ so $\nabla\!\cdot\!\mathbf{B}=0$) and scanned $B_\phi/B_z=0$--$2$ ($B_\phi=0$ reproduces the longitudinal spectrum, a code check). The axial growth-rate plateau peaks at $\lambda\approx1\,\lambda_{\rm J}$ for every toroidal fraction, with no excess growth at shorter $\lambda$ even for a strongly-wound field (Fig.~\ref{fig:helical_scan}a). The toroidal component's effect is a radial \emph{pinch} narrowing the filament ($W_{\rm FWHM}$: $0.57\to0.43\,\lambda_{\rm J}$, $\approx25\%$; Fig.~\ref{fig:helical_scan}b), so it raises $\lambda/W$ by reducing $W$ --- moving it \emph{towards} the classical value. Crucially this field is \emph{seeded}, not in force balance: imposed on an already-contracting filament, it drives an out-of-equilibrium pinch rather than the confining toroidal equilibrium whose sausage ($m=0$) mode carries the shortened \citet{Fiege2000} wavelength, so it \emph{cannot} test that mechanism. The static longitudinal and perpendicular geometries are thus disposed of as sources of a short axial spacing (longitudinal leaves $\lambda_{\rm max}$ unchanged; perpendicular spindles), but the helical/\citet{Fiege2000} pinch --- the one geometry that could shorten $\lambda$ --- remains \emph{untested}: it requires a force-balanced helical equilibrium, deferred to future work. Magnetic tension is thus not the demonstrated cause, but neither is it excluded in its one relevant configuration.

\begin{figure}
\includegraphics[width=0.48\textwidth]{figures/fig_perp_ladder.pdf}
\caption{Ideal perpendicular-field refinement ladder ($128^3$--$512^3$; Section~\ref{sec:optional_tests}). \emph{(a)} Cells per local Jeans length $N_J$ at the beading-relevant spine contrast $C\approx3$: the collapsing spine is resolved far above the Truelove threshold ($N_J=4$, dashed) at every resolution. \emph{(b)} The spurious short-$\lambda$ axial ``beading'' emerges only once the spine drops below $N_J\approx1$, and refinement pushes that onset to ever-deeper collapse (at $512^3$ the filament stays monolithic to $N_J\approx40$ at $C\approx30$) --- the signature of a grid artefact, not a physical fragmentation mode.}
\label{fig:perp_ladder}
\end{figure}

\begin{figure}
\includegraphics[width=0.48\textwidth]{figures/fig_helical_scan.pdf}
\caption{Divergence-free helical-field scan ($B_\phi/B_z=0$--$2$, $512\times128^2$; Section~\ref{sec:optional_tests}). \emph{(a)} The axial growth-rate spectrum $\Gamma(\lambda)$ in the physical band is invariant to the toroidal fraction ($<6\%$, fit noise), with the plateau peak fixed at $\lambda\approx1\,\lambda_{\rm J}$; $B_\phi/B_z=0$ (black) reproduces the pure-longitudinal spectrum, a code check. No excess growth appears at shorter $\lambda$ for any winding. \emph{(b)} The one systematic effect is a radial \emph{pinch}: the formation FWHM narrows by $\approx25\%$, with onset at $B_\phi\gtrsim B_z$. The seeded helical field therefore raises $\lambda/W$ by shrinking $W$, not by shortening $\lambda$; the force-balanced \citet{Fiege2000} pinch mode is not tested here.}
\label{fig:helical_scan}
\end{figure}

\section{The fragmentation scale and its comparison with observation}
\label{sec:validation_endtoend}

The measured core spacing is sub-classical only when expressed against the filament FWHM, and that comparison hinges on how the width is defined. We therefore anchor it on the spacing in units of the filament FWHM measured the same way on simulations and data, and verify the whole chain end to end (Appendix~\ref{sec:val_appendix}).

First, the measurement is unbiased: passing synthetic filaments with a known spacing and width through the identical pipeline recovers $\lambda/W$ to better than $0.5\%$ over $\lambda/W=1.4$--$5.7$ (Appendix~\ref{sec:val_appendix}), so the observed$\to$simulated comparison depends on the width convention, not on measurement bias. Second, forward-modelling the simulations end to end through the HGBS pipeline (projection, beam convolution, Plummer fitting, and the same NN measurement) gives the observable $\lambda/W_{\rm fil}$ directly. The forward-model value ($\approx1.0$--$1.3$) is $\approx30\%$ above the analytic shortcut ($\approx0.9$, Eq.~\ref{eq:t1_apply}), which applies only a single $\sigma\!\to\!$FWHM factor; we adopt the forward-model value as definitive and carry this $\approx30\%$ offset as the width-conversion systematic.

The fragmentation wavelength must be measured with care: counting density peaks late in the collapse is epoch- and resolution-dependent, so we instead identify the dominant mode from the growth-rate spectrum. Seeding the filament with a small broadband perturbation ($\delta v/c_s=10^{-3}$, twelve axial modes), we track the power $P(k,t)$ in each mode and fit its growth over its approximately exponential interval (Fig.~\ref{fig:growthcurves}); the growth accelerates as the filament collapses, so we isolate the dominant mode of a \emph{dynamically collapsing} filament. The growth rate $\Gamma(\lambda)$ (Fig.~\ref{fig:growthrate}, left) rises as $\lambda$ decreases and is resolution-stable for $\lambda\gtrsim1\,\lambda_{\rm J}$ (low-$n$ modes agree to $\sim5\%$ across $128^3$--$512^3$; $\Gamma$ at $n=7$ is $17.8,18.5,18.7$), rising to a \emph{plateau} at the converged-band edge ($\Gamma$ at $n=6,7,8$, i.e.\ $\lambda=1.33,1.14,1.00\,\lambda_{\rm J}$, is $17.6,18.7,18.7$), while below $\lambda\lesssim0.9\,\lambda_{\rm J}$ it rises with \emph{transverse} resolution and is numerical. Because $\Gamma$ plateaus rather than turning over, we quote the dominant wavelength as an \emph{upper bound}, $\lambda_{\rm max}\lesssim1\,\lambda_{\rm J}$, and do \emph{not} claim a precisely-located maximum (Appendix~\ref{sec:val_appendix}); the true peak could lie at shorter, transverse-unresolved wavelength. It is an order-$\lambda_{\rm J}$, sub-classical value: $\lambda_{\rm max}/W_{\rm FWHM}\approx1.5$--$2.0$, a factor $\approx2.5$--$3$ below the classical $\lambda_{\rm m}\approx5$--$6\,W_{\rm FWHM}$ (\citealt{Fischera2012}) and at or below the critical $2.6\,W_{\rm FWHM}$ across $f=1.1$--$2.0$. The observations use a beam-convolved Plummer width, so the \emph{like-for-like} value (below) is $\lambda/W\approx1.0$--$1.3$, not $\approx2$, and it is this matched value that should be set against the data (Table~\ref{tab:wave}).\footnote{The near-critical formation FWHM is $W_{\rm FWHM}=0.58\,\lambda_{\rm J}$, narrower than the supercritical $W_{\rm form}=0.683\,\lambda_{\rm J}$ because a near-critical filament has already begun to contract when its mode is set (both tabulated in Table~\ref{tab:widthdefs}).}

\begin{figure}
\includegraphics[width=0.47\textwidth]{figures/fig_growthcurves.pdf}
\caption{Power in three axial modes versus time for a near-critical filament ($f=1.1$). Each grows with an approximately exponential interval (the fit window, shaded) whose slope gives the growth rate; the growth steepens as the filament collapses, so the extracted rate characterises the dominant mode of a dynamically collapsing filament. The peak mode $n=7$ grows faster than the long-wavelength $n=3$, giving the sub-classical $\lambda_{\rm max}$; the grid-scale $n=12$ is shown for reference.}
\label{fig:growthcurves}
\end{figure}

\begin{figure*}
\includegraphics[width=0.80\textwidth]{figures/fig_growthrate.pdf}
\caption{\emph{(a)} Growth-rate spectrum $\Gamma(\lambda)$ of the axial modes for a near-critical filament ($f=1.1$) at three resolutions. For $\lambda\gtrsim1\,\lambda_{\rm J}$ the curves overlap (resolution-converged); at the grid scale ($\lambda\lesssim0.9\,\lambda_{\rm J}$, shaded) they separate and rise with resolution, the signature of unresolved sub-fragmentation. $\Gamma$ plateaus (does not turn over) at the converged-band edge, so $\lambda_{\rm max}\lesssim1\,\lambda_{\rm J}$ is an \emph{upper bound}, not a resolved maximum. \emph{(b)} $\lambda_{\rm max}/W_{\rm FWHM}$ versus line-mass fraction (Gaussian FWHM), well below the classical $5$--$6$ at every $f$. The like-for-like comparison uses a beam-convolved Plummer width $\sim1.5\times$ broader, giving matched $\lambda/W\approx1.0$--$1.3$ (grey band): this, not the Gaussian $\approx2$, is what should be set against the observed $\lambda/W\approx0.4$--$1.9$.}
\label{fig:growthrate}
\end{figure*}

The convention-independent statement is that the dominant fragmentation wavelength of a dynamically collapsing filament is sub-classical. In width-free form $\lambda_{\rm max}\lesssim1.0$--$1.14\,\lambda_{\rm J}$ versus the classical $\lambda_{\rm m}=22H=2.3\,\lambda_{\rm J}$ (a factor $\approx2$), which we take as primary; in FWHM units $\lambda_{\rm max}/W_{\rm FWHM}\approx2$ versus the classical $5$--$6$ (a factor $\approx2.5$--$3$, larger only because the contracting-run FWHM $0.58\,\lambda_{\rm J}$ exceeds the equilibrium $\approx0.45\,\lambda_{\rm J}$ on which $5$--$6$ is defined), at or below even the critical $2.6\,W_{\rm FWHM}$. This is an \emph{upper bound}: the spectrum plateaus at the converged-band edge ($\lambda\approx1\,\lambda_{\rm J}$) and does not turn over, so the true peak may lie shorter ($\lambda_{\rm max}/\lambda_{\rm m}\approx0.35$--$0.5$ across $f=1.1$--$2.0$, each an upper bound). The selected \emph{mode} is imprinted early and fit-window-invariant (Appendix~\ref{sec:val_appendix}), by which epoch the filament has contracted to $C\approx20$--$40$ with the Jeans length still resolved ($N_J\gtrsim5$). The equilibrium-recovery control (Appendix~\ref{sec:val_appendix}) returns the classical $5$--$6$ for a true equilibrium and $\approx1.4$--$2$ for a contracting one, so the sub-classical ratio is physical, in the direction predicted for an accreting/contracting filament \citep{Clarke2016,Clarke2017}.

Two caveats temper the causal attribution. First, the fitted rates ($\Gamma\approx18$--$26\,t_{\rm J}^{-1}$) are not quasi-static eigenvalues but the accelerating growth of a collapsing background (they match the free-fall rate at $\rho_{\rm eff}/\rho_0\approx8$), so we use only the mode \emph{ordering}, not the absolute rate. The selected wavelength is not the instantaneous Jeans scale of the contracted core --- $\lambda_{\rm max}\approx1\,\lambda_J(\rho_0)\approx2.8\,\lambda_J(\rho_{\rm eff})$, nearly three times it --- but is close to the \emph{classical} scale evaluated at the elevated density, $22H(\rho_{\rm eff})\approx0.8\,\lambda_J(\rho_0)$: the filament fragments at approximately the classical scale \emph{for its current density}. The sub-classical \emph{deficit} in $\lambda/W$ thus arises not from non-classical physics but because the observable FWHM $W$ does not contract in step with $\rho_c$ (the extended envelope lags the collapsing spine), so $\lambda\sim22H(\rho_c)$ shrinks while $W$ stays broad; the equilibrium control, in which $W$ and $22H$ are locked, confirms this by returning $5$--$6$. Second, our Plummer-like ($p\simeq2$) profile matches observed filaments but not the isothermal equilibrium ($p=4$); a $p=4$ profile gives the same order-$\lambda_{\rm J}$ wavelength ($\approx0.8\,\lambda_{\rm J}$), so the deficit is not a profile-index artefact, but that cylinder too is only approximately relaxed, so what we demonstrate is generic ``fragmentation from a contracting state'' rather than specifically the accretion-driven mechanism (Section~\ref{sec:limitations}).

The simulated observable is worth converting once. The internal ratio is $\lambda_{\rm max}/W_{\rm core}\approx3.3$; the $\sigma$-normalised shortcut ($\times0.27$; Eq.~\ref{eq:t1_apply}) gives $\lambda/W_{\rm fil}\approx0.9$, whereas forward-modelling the same runs through the HGBS pipeline gives $\lambda/W_{\rm fil}\approx1.0$--$1.3$, which we adopt as \emph{the} observable (the shortcut normalises by the initial $\sigma$, not the evolved Plummer FWHM). In Gaussian formation-FWHM units the same wavelength is $\approx1.5$--$2.0$, but a Gaussian FWHM is $\sim1.5$--$1.9\times$ narrower than the observed Plummer FWHM, so the ``$\approx2\approx1.9$'' coincidence is not like-for-like. This comparison is resolution-limited as well as convention-limited: the simulated wavelength is an upper limit that could shift to shorter $\lambda$ at higher resolution, moving single-filament fragmentation from sufficiency toward insufficiency.

Single-filament gravitational fragmentation, evolved dynamically, thus produces a converged \emph{upper bound} on the fragmentation wavelength that is sub-classical and consistent with the observed spacing for the near-critical filaments that dominate the sample. What is bounded is the wavelength ($\lambda_{\rm max}\lesssim1\,\lambda_{\rm J}$), not a located maximum: the spectrum does not turn over within the converged band, and the object measured is the dominant mode of a \emph{dynamically collapsing} configuration, so the robust deliverable is the density-independent ratio $\lambda/W$ with its equilibrium-recovery calibration, not a precise linear $\lambda_{\rm max}$. That control is the load-bearing evidence: a pressure-confined Ostriker cylinder is drag-relaxed to the solver's hydrostatic state (a no-perturbation run then drifts only $\sim3\%$ in central density, in the $k=0$ mode), the drag removed, and the \emph{same} seed and fit applied; the identical pipeline recovers the classical turnover at $22H$ ($\lambda_{\rm max}/22H=1.07$--$1.10$, both resolutions; the $\approx7$--$10\%$ long bias is conservative), whereas the contracting runs have no turnover down to the grid-limited band (full method in Appendix~\ref{sec:val_appendix}). This is a numerical instance of the generic expectation that a still-contracting filament fragments below the marginally-stable hydrostatic mode \citep{Clarke2016,Clarke2017,Clarke2020}, not a new mechanism; our distinct contributions are the longitudinal near-critical demonstration, the equilibrium-recovery control, and the homogeneous Gaia-DR3 reanalysis.

Because the simulated upper bound ($\lambda/W\approx1.0$--$1.3$) already sits at or just below the adopted observed $(\lambda/W)_0\approx1.4$, single-filament dynamical fragmentation is \emph{at best marginally} sufficient and cannot be claimed as the sole mechanism. And because the simulated scale is only an upper bound, this establishes consistency, not sufficiency: a shorter true wavelength would make single-filament fragmentation \emph{over}-produce cores, so neither uniqueness nor sufficiency is demonstrated, and hierarchical-fibre projection (Section~\ref{sec:fibre_projection}) and magnetic tension remain viable within the factor-$\sim$2 systematics. The claim is falsifiable in two operational senses. It is \emph{refuted} if, on resolved single filaments or individual fibres measured at a fixed width convention, the core line-density is no higher than the classical one-core-per-$\lambda_{\rm m}$ (cores as sparse as $\lambda_{\rm m}\approx5$--$6\times$FWHM). Conversely, given the simulated upper bound $\lambda/W\approx1.0$--$1.3$, if the true observed $\lambda/W$ instead fell well below $\sim1$ single-filament fragmentation would \emph{over}-produce cores and be disfavoured, so the mechanism is operationally viable only in the band between these limits. Fibre-resolved kinematics placing cores at the classical spacing \emph{within} individual fibres would favour projection, and more strongly supercritical filaments should show shorter core spacing (Section~\ref{sec:limitations}).

\emph{Opacity-limited fragmentation.} Because $\Gamma(\lambda)$ does not turn over within the converged band, we cannot exclude that it keeps rising until non-isothermal physics intervenes at a scale below anything modelled here. The isothermal runs contain no small-scale stabilising physics, so the reported ceiling $\lambda_{\rm max}\lesssim1\,\lambda_{\rm J}$ is contingent on isothermality holding to that scale: it is a numerical ceiling set by the resolved band, not a guaranteed physical minimum fragment scale. In reality that minimum is set by the opacity limit, where the gas becomes optically thick to its cooling radiation and an effective equation-of-state transition halts fragmentation; our EOS tests (Appendix~\ref{sec:validation}) show that a $\gamma=5/3$ (adiabatic) equation of state suppresses fragmentation, so the true short-wavelength cutoff lies below our isothermal grid and is not captured here.

\section{DISCUSSION}
\label{sec:discussion}

\subsection{Why the spacing matters: fragment masses and the CMF}
\label{sec:cmf}
The spacing matters because, with the line mass, it sets the characteristic fragment mass that seeds the core mass function (CMF) and ultimately the IMF. For a filament of line mass $M_{\rm line}$ fragmenting at spacing $\lambda$, the mean fragment mass is $M_{\rm frag}\sim\epsilon\,M_{\rm line}\,\lambda$ ($\epsilon\lesssim1$ the reservoir fraction). A near-critical $M_{\rm line}\approx16\,M_\odot\,{\rm pc^{-1}}$ and $\lambda\approx0.1$--$0.2$\,pc give $M_{\rm frag}\sim1.6$--$3\,M_\odot$, i.e.\ $\sim0.5$--$1\,M_\odot$ for $\epsilon\sim0.3$, consistent with the HGBS CMF peak; the classical $\approx5$--$6\times$FWHM would place cores $3$--$4\times$ farther apart and overshoot it. This is an order-of-magnitude \emph{consistency check}, not a fit: $\epsilon$ and $M_{\rm line}$ are degenerate (only $\epsilon M_{\rm line}$ enters), and because the observable is a line density (Section~\ref{sec:clustering}) it predicts a \emph{mean} fragment mass, not a peaked CMF. The robust point is the sign, not a numerical CMF prediction.


\subsection{Hierarchical-fibre projection is consistent with the sub-classical spacing}
\label{sec:fibre_projection}

A second, purely geometric route reaches the same sub-classical value without modified physics, and rests on independently observed substructure: HGBS filaments are bundles of $N$ velocity-coherent \emph{fibres} \citep{Hacar2013,Hacar2018,Yang2024}. If each fibre fragments near the classical wavelength but core spacing is measured \emph{across} the bundle, as any filament-level pipeline does, the interleaving of cores from $N$ fibres compresses the apparent spacing. The premise --- quasi-periodic fibre-to-core fragmentation at the classical scale --- is a hypothesis, not an established measurement (\citealt{Yang2024} find only a weak imprint), so we present projection as viable but unconfirmed.

We tested this with a forward-model that reuses our exact NN measurement (order cores along the spine by arclength; take adjacent separations; median $=\lambda$). Each fibre is a cylinder of width $W_{\rm fib}$ fragmenting at its \emph{own} classical wavelength $\lambda_{\rm fib}\approx5$--$6\,W_{\rm fib}$ (\citealt{Fischera2012}); we populate a filament with $N$ such fibres (random phase, $\sim30\%$ lognormal scatter), project onto the shared spine, and measure in units of the \emph{filament} width $W_{\rm fil}$. This assumes fibres are quasi-hydrostatic, which sits uneasily with our own result: a subcritical fibre would not fragment at all, whereas a near-critical one would by our argument fragment sub-classically at $\approx\tfrac12\lambda_{\rm fib}$, doubling the $r/N$ needed to match the data (Eq.~\ref{eq:fourrn} becomes $(2.5$--$3)r/N$). We therefore treat the classical-fibre $(5$--$6)r/N$ result as an upper bound; the fibre and single-filament pictures are not independent. Two effects compress the spacing --- narrower fibres ($r\equiv W_{\rm fib}/W_{\rm fil}<1$) and several of them --- giving in the interleaving limit
\begin{equation}
\frac{\lambda}{W_{\rm fil}} \;\approx\; \frac{(5\text{--}6)\,r}{N},
\label{eq:fourrn}
\end{equation}
verified across $r$ and $N$ in Fig.~\ref{fig:fibre_degen}. Either effect alone can produce a sub-classical value: a single fibre of half the filament width ($r=0.5$, $N=1$) gives $\lambda/W_{\rm fil}=2.0$, as do two full-width fibres ($r=1$, $N=2$).

The observation constrains only $r/N$, not $N$ individually. The fixed-width value ($\lambda/W_{\rm fil}\approx1.9$) requires $r/N\approx0.35$; the region-Plummer value ($\approx1.4$) gives $r/N\approx0.25$; FilFinder ($0.4$--$0.8$) gives $r/N\approx0.07$--$0.15$. Observed fibre widths $W_{\rm fib}\approx0.03$--$0.05$\,pc ($r\approx0.3$--$0.5$; \citealt{Hacar2013,Hacar2018}) then imply $N\approx1$--$2$ for the DisPerSE spacings and $N\approx3$--$5$ for FilFinder. \citet{Hacar2013} resolve $3$--$5$ fibres per bundle in Taurus, giving $r/N\approx0.07$--$0.15$, below the $\approx0.35$ the DisPerSE spacing requires --- so at face value the projection model would \emph{over}-compress, matching DisPerSE only at $N\approx1$--$2$, and is internally consistent only in a restricted $(r,N)$ corner. We do not claim a specific fibre number; the robust statement is that measuring across narrower, multiple fibres reduces the apparent spacing below the classical prediction with no modified physics.

The $(5$--$6)r/N$ scaling is also reproduced when the forward-model is run on the real DisPerSE arclengths and counts of the four robust regions rather than idealised straight spines; that exercise is only a consistency check, since it seeds classical-spacing cores onto the observed skeletons and so assumes the fibre population rather than measuring it.

\begin{figure}
\includegraphics[width=0.48\textwidth]{figures/fig_fibre_degeneracy.pdf}
\caption{Hierarchical-fibre forward-model, showing the filament-level $\lambda/W_{\rm fil}$ recovered by our NN pipeline as a function of the fibre-to-filament width ratio $r=W_{\rm fib}/W_{\rm fil}$ and the number of fibres $N$; the model follows $\lambda/W_{\rm fil}\approx(5$--$6)r/N$ (Eq.~\ref{eq:fourrn}). White contours mark the observed fixed-width and region-specific values, the orange contour the FilFinder value, and the red band the observed fibre-width range ($r\approx0.3$--$0.5$). The observation constrains the combination $r/N$, not $N$ alone: the sub-classical value is reproduced by a single narrow fibre, by a few wider fibres, or by intermediate combinations.}
\label{fig:fibre_degen}
\end{figure}

Hierarchical projection is thus a natural and sufficient way to produce the sub-classical spacing without modified magnetic physics, and it also accounts for the skeleton-algorithm and $L/3$ systematics (Section~\ref{sec:error_budget}). It is not \emph{required}: single-filament dynamical fragmentation already accounts for the bulk of the spacing (Section~\ref{sec:validation_endtoend}), so projection is one of at least two viable readings, and it assumes rather than measures the fibre population. The decisive test is to decompose the observed cores into their parent fibres and measure the fibre-level spacing across the robust regions (Section~\ref{sec:limitations}). Consistent with this under-determination, \citet{Clarke2020} find that fragmentation through sub-filaments leaves \emph{no} single characteristic length-scale.

A coherence point remains: our single-filament simulations fragment into a \emph{quasi-periodic} comb at $\lambda_{\rm max}$, whereas the observed cores are clustered with no comb (Section~\ref{sec:clustering}). The two reconcile only if sub-fragmentation, multi-fibre interleaving, and projection erase the periodicity into the observed near-exponential gap distribution, as \citet{Clarke2020} find. We do not simulate this stage, so the match between a single-filament $\lambda_{\rm max}$ and the observed line density is an order-of-magnitude, not a morphological, agreement.



\subsection{Field geometry, boundary conditions and the perpendicular case}
\label{sec:why_mhd_fails}

Across the campaign the \emph{outcome} of a run --- radial collapse or longitudinal beading --- depends on regime and boundary conditions, so a single filament has no setup-independent fate; its dominant fragmentation \emph{mode} is nonetheless resolution-stable (Section~\ref{sec:validation_endtoend}), a converged \emph{upper bound} $\lesssim0.5$ times the classical value. It is sometimes argued that most filaments are threaded by perpendicular fields, from the tendency of dense filaments to lie perpendicular to the large-scale Planck field \citep{Planck2016}; but the projected orientation need not equal the internal three-dimensional geometry, and hourglass morphologies in which the field turns longitudinal inside the crest are observed \citep{Jadhav2026}. A resolved perpendicular filament spindles rather than fragmenting axially (Section~\ref{sec:optional_tests}), so it is not a source of tension: perpendicular filaments do not produce the observed axial core chains.

Where perpendicular geometry does apply, three non-exclusive resolutions of the residual deficit remain: hierarchical fibre fragmentation (bundles of near-critical fibres, each fragmenting longitudinally and projected across the bundle; \citealt{Hacar2013,Hacar2018,Yang2024}); non-ideal MHD at later evolution, which narrows the factor-$\sim3$ perpendicular deficit in longer runs; and an observational lag, in which the observed cores are relics of an earlier fragmentation episode. The two leading interpretations, single-filament dynamical fragmentation and hierarchical projection, are distinguished decisively by fibre-resolved core spacing across the robust regions and by high-resolution polarimetry of filament interiors (Section~\ref{sec:limitations}); the temperature-gradient signature (Section~\ref{sec:mt_variation}) adds a third test, predicting spatially differential fragmentation that fibre geometry alone would not produce.

Sub-isothermal equations of state ($\gamma<1$) accelerate fragmentation by $30$--$40\%$ without changing the wavelength, while adiabatic heating ($\gamma=5/3$) suppresses it entirely (Appendix~\ref{sec:validation}), so the isothermal suite is the conservative case in a globally averaged sense --- though not locally, where a $\sim12\%$ dust-temperature variation shifts the local $\lambda/W$ by $\sim10$\,per-cent and relocates where cores form (Section~\ref{sec:mt_variation}).

\subsubsection{Non-Ideal MHD Effects}
\label{sec:nonideal}

At the densities probed here ($n_0\sim10^4$\,cm$^{-3}$), the ambipolar-diffusion time ($t_{\rm AD}\propto L^2/\eta_{\rm AD}$; $\sim10$--$20\,t_{\rm ff}$ for typical filament conditions) is long compared with the collapse time \citep{Tassis2004}, so ideal MHD is a good approximation for the longitudinal fragmentation stage: four non-ideal runs at $\eta_{\rm AD}=\{0.01,0.05\}$ still collapse radially, $\sim12$--$15\%$ \textit{earlier} than the ideal control (flux leakage erodes magnetic support, the conservative direction). Non-ideal effects nonetheless change the outcome for perpendicular fields, where ambipolar diffusion promotes longitudinal fragmentation (Section~\ref{sec:modecompetition}); Hall and Ohmic terms remain relevant for late-stage core evolution.

\subsection{Limitations and future work}
\label{sec:limitations}

Two limitations bound our conclusion. First, the characteristic $\lambda/W$ and its scatter ($\sigma\approx0.4$) rest on only four nearby, low-mass Gould Belt regions of broadly solar metallicity and similar radiation field, so the result is established only for that regime; whether it persists in high-mass complexes, at lower metallicity, or under stronger irradiation is untested. Four clouds are a thin base for an intrinsic-scatter estimate or for claiming the adopted $(\lambda/W)_0\approx1.4$ is representative of filaments generally, and the value carries the factor-of-two skeleton-algorithm and width-convention systematics (Section~\ref{sec:error_budget}); extending the pipeline to further HGBS regions and independent surveys (e.g.\ JCMT/SCUBA-2) is needed. Two validation checks also remain \emph{outstanding}: the $\lesssim5\%$ manual-editing bound is a \emph{proxy} from the junction/association tests pending a blind, multi-editor re-skeletonisation, and a dedicated faint-segment injection--recovery test is not yet done; neither should be read as complete. Second, the simulations are idealised controlled experiments, not models of specific filaments: the periodic-BC main grid does not represent the turbulent, magnetically coupled boundaries of a real filament, and our ambient-confined configuration is an idealisation of that confinement.

On the causal attribution: our seeded filaments contract from the start, so what we demonstrate is generic ``fragmentation from a non-equilibrium, contracting state'' rather than specifically the accretion-driven case of \citet{Clarke2016,Clarke2017}; a near-equilibrium filament grown by controlled accretion is the decisive future test. The equilibrium-recovery control itself is a \emph{two-point} comparison (one relaxed equilibrium versus one contracting run): it establishes the \emph{sign} of the contraction effect but not its functional dependence on the degree of non-equilibrium, for which a relaxation ladder of partial-relaxation states is the needed future test to make the causal link quantitative. Further idealisations are: (i) an \emph{isothermal} EOS (locally uncertain at $\sim10$--$20\%$; Section~\ref{sec:mt_variation}); (ii) \emph{periodic transverse boundaries}, whose accretion channel suppresses fragmentation unless ambient confinement is imposed (Section~\ref{sec:modecompetition}); (iii) \emph{simplified initial profiles}; and (iv) a \emph{prescribed Kolmogorov turbulence spectrum} rather than driven transonic turbulence, leaving the wavelength unchanged ($<5\%$) but shortening the timescale $3$--$5\times$ (Appendix~\ref{sec:supplementary_campaigns}).

The fragmentation \emph{classification} is resolution-independent (the collapse \emph{time} itself carries an $8.5\%$ offset between $128^3$ and $256^3$ and is not claimed converged; Appendix~\ref{sec:validation}) and domain-doubling changes the collapse time by $<2\%$, but the transition zone is sampled with only 2 seeds per point. Non-linear core merging could in principle inflate the apparent supercritical spacing by up to $\sim50\%$ over a few Myr, but we have verified it does \emph{not} bias the quoted value: $\lambda/W_{\rm fil}$ is measured at maximum bead count (fragmentation onset), and in $38/39$ runs the peak count is still rising at the final snapshot while no run shows the peak-count decline that would signal merging. The ambipolar wavelengths are likewise fragmentation-onset measurements (the runs reach gravitational runaway before a converged plateau is established; Section~\ref{sec:modecompetition}).

The decisive follow-ups are: (i) fibre-resolved NN spacing on the full HGBS catalogues, to test the hierarchical interpretation; (ii) high-resolution polarimetry of filament interiors, to test the longitudinal-versus-perpendicular assumption; (iii) MHD simulations with a spatially varying temperature field (Section~\ref{sec:mt_variation}); (iv) longer non-ideal runs, to test whether the perpendicular deficit closes; (v) a seeding-fixed domain study to \emph{refine} (rather than bound) the near-critical $\lambda_{\rm max}$, the only remaining lever on its precision (the seed-realisation and open-gravity-boundary robustness having been established, Appendix~\ref{sec:val_appendix}).


\subsection{Non-isothermal effects and the temperature-gradient test}
\label{sec:mt_variation}

The MHD suite is isothermal, whereas HGBS dust-temperature maps vary by $\sim12$\,per-cent ($\sim10$--$13$\,K), a local perturbation since $\lambda_{\rm J}\propto\sqrt{T}$. A single illustrative run with an imposed $\sim12$\,per-cent longitudinal temperature gradient fragments \emph{differentially} (the cool, locally supercritical half beads while the warm half is suppressed), shifting the local wavelength by $\sim10$\,per-cent, because temperature sets both the Jeans scale and the critical line mass ($\mu_{\rm crit}\propto T$). The isothermal assumption is thus conservative only in a globally averaged sense; a self-consistent treatment is deferred to future work.


\section{CONCLUSIONS}

We have investigated the origin of the sub-classical core spacing in HGBS filaments, combining a Gaia DR3 observational re-analysis, a large MHD campaign, and a hierarchical-fibre forward-model. The quoted $(\lambda/W)_0\approx1.4$ is a descriptive characterisation of four clouds, not a population estimate with a formal confidence level. Our primary result is \emph{under-determination}: the spacing is sub-classical in \emph{sign} under every convention, but its magnitude is convention-dominated (the skeleton algorithm alone moves $\lambda/W$ by a factor of two) and we do not identify the cause. The simulations demonstrate \emph{plausibility} --- a contracting filament \emph{can} fragment sub-classically, a resolution-stable upper bound whose spectrum does not turn over --- not that they explain the observed spacing. The deficit is consistent with dynamic collapse, but the force-balanced helical/toroidal magnetic pathway remains an \emph{un-eliminated} alternative. We group the conclusions as follows.

\begin{itemize}
    \item \textbf{Robust observational findings.}
    \begin{itemize}
        \item The nearest-neighbour spacing is sub-classical in \emph{sign} under every skeleton and width convention, by a factor $\approx3$--$4$ (DisPerSE; more under FilFinder) relative to the classical $\lambda_{\rm m}\approx5$--$6\times$FWHM.
        \item Under directly-fitted region-Plummer widths the adopted spacing is $(\lambda/W)_0\approx1.4$ (fixed-width $\approx1.9$--$2.1$; all sub-classical), robust in sign against population splits, bias correction, hierarchical bootstrap and an independent FilFinder cross-check.
        \item The cores reject \emph{both} periodic and pure-Poisson spacing, consistent with a clustered line density bearing a beam-scale hard-core exclusion rather than a resonant wavelength.
    \end{itemize}

    \item \textbf{Methodological findings.}
    \begin{itemize}
        \item The pairwise median is a biased estimator that converges to the region extent $L/3$, not the fragmentation wavelength; the nearest-neighbour spacing is the appropriate observable.
        \item The magnitude of $\lambda/W$ is convention-dominated, the skeleton algorithm (DisPerSE vs FilFinder, a factor of two) being the single largest term and the dominant limit on any comparison with theory.
        \item The measurement pipeline is unbiased: applied to a relaxed force-balanced equilibrium cylinder it recovers the classical $\lambda_{\rm m}=22H\approx5$--$6\,W_{\rm FWHM}$, whereas the \emph{same} pipeline gives $\approx1.4\,W_{\rm FWHM}$ for contracting filaments.
    \end{itemize}

    \item \textbf{Simulation findings (plausibility, not proof).}
    \begin{itemize}
        \item A dynamically contracting filament \emph{can} fragment sub-classically at a resolution-stable \emph{upper bound} ($\lambda_{\rm max}\lesssim1\,\lambda_{\rm J}$; the spectrum does not turn over), demonstrating plausibility but neither uniqueness nor sufficiency.
        \item In matched Plummer units the simulated $\lambda/W\approx1.0$--$1.3$ matches the region-Plummer observations, under-predicts fixed-width DisPerSE and over-predicts FilFinder, so whether single-filament fragmentation suffices depends on the segmentation convention.
        \item Hierarchical-fibre projection reaches the same value geometrically ($\lambda/W_{\rm fil}\approx(5$--$6)r/N$) and remains a viable, but not required, explanation.
    \end{itemize}

    \item \textbf{Unresolved physics.}
    \begin{itemize}
        \item We do not identify the cause: single-filament dynamical fragmentation, hierarchical fibres and (helical) magnetic tension all remain open, and current data cannot discriminate them.
        \item An \emph{ideal-MHD} perpendicular field, resolved to $N_J\approx40$--$170$, drives monolithic radial collapse with no axial fragmentation (the short low-resolution beading is a Truelove artefact), so a static perpendicular field cannot be the source; a \emph{seeded} helical field does not shorten $\lambda_{\rm max}$, so the force-balanced helical/toroidal pinch remains \emph{untested} and un-eliminated.
        \item A single illustrative run indicates the isothermal assumption can be locally non-conservative under a $\sim12\%$ temperature gradient; a fuller treatment, and the outstanding blind re-skeletonisation check, are deferred.
        \item The decisive test is fibre-resolved NN spacing on the full HGBS catalogues, with high-resolution polarimetry of filament interiors, a perpendicular-field resolution study, and MHD runs with a spatially varying temperature field as complementary probes.
    \end{itemize}
    \end{itemize}

\section*{ACKNOWLEDGMENTS}

We thank the Herschel Gould Belt Survey team for the datasets. HGBS data are available from the Herschel Science Archive. The Athena++ simulations were run on cloud computing infrastructure (Google Cloud E2); the full set of $2251$ Athena++ MHD simulations (Appendix~\ref{sec:campaigns}) consumed of order $10^5$ CPU-hours.

\section*{DATA AVAILABILITY}

A public Zenodo archive accompanies this paper (DOI reserved; deposited at submission for referee access) and contains: (i) the full run manifest (all 2251 runs with their $(f,\beta,\mathcal{M},\theta,\eta_{\rm AD})$ parameters and outcomes); (ii) the Athena++ problem-generator source (\texttt{filament\_ambient.cpp}, including the divergence-free helical construction) and normalisation scripts; (iii) HDF5 snapshots of the growth-rate convergence subset ($128^3$--$512^3$) and the equilibrium-recovery control (remaining snapshots on request); and (iv) the analysis pipeline (DisPerSE v0.9.24 parameters, the bridging/deblending and nearest-neighbour scripts, and the derived $\lambda/W$ catalogues). The estimator was independently re-implemented within our group and checked against the deposited counts and medians (Table~\ref{tab:nnpipeline}); to date the nearest-neighbour and skeleton pipelines have \emph{not} been independently re-run by anyone other than the author, and we release the code and catalogues to enable that. HGBS column-density and skeleton maps are available from the Herschel Science Archive.

\bibliographystyle{mnras}
\appendix

\section{Validation of the measurement and the width normalisation}
\label{sec:val_appendix}

This appendix documents the three checks summarised in Section~\ref{sec:validation_endtoend}.

To verify that the routines extracting $\lambda$ and $W$ carry no hidden normalisation bias, we built synthetic filament cubes at the simulation resolution with a \emph{known} transverse FWHM $W_{\rm in}$ and a \emph{known} axial bead spacing $\lambda_{\rm in}$ (localised Gaussian beads on a Gaussian transverse profile) and passed them through the identical pipeline used on the simulations (spine-averaged transverse profile, Gaussian and beam-convolved Plummer width fits, axial peak-spacing). The recovered $\lambda/W$ matches the input to better than $0.5\%$ over $\lambda/W_{\rm in}=1.4$--$5.7$ (Fig.~\ref{fig:pipeval}). The measurement is therefore unbiased, and the corrected width bookkeeping of Eq.~(\ref{eq:t1_apply}) is confirmed.

\begin{figure}
\includegraphics[width=0.46\textwidth]{figures/fig_pipeline_validation.pdf}
\caption{Closed-loop validation: the pipeline recovers the injected $\lambda/W$ to within $0.4\%$ over $\lambda/W=1.4$--$5.7$.}
\label{fig:pipeval}
\end{figure}

\emph{Growth-rate measurement and convergence.} We take the fragmentation wavelength as the peak of the growth-rate spectrum, not a late-time peak count (which sub-fragments with resolution). Seeding eight ambient-confined filaments ($f=1.1,1.2,1.5$; $\beta=1$; two seeds) with a broadband perturbation ($\delta v/c_s=10^{-3}$, twelve axial modes) and fitting each mode's growth over $\delta_{\rm rms}=3\times$seed$\to0.3$ gives $\Gamma(\lambda)$ (Fig.~\ref{fig:growthcurves}). The spectrum is resolution-converged for $\lambda\gtrsim1\,\lambda_{\rm J}$: for $f=1.1$ the $n=7$ mode has $\Gamma=17.8,18.5,18.7$ at $128^3,256^3,512^3$, flat through the converged-band edge ($\Gamma(n{=}6,7,8)=17.6,18.7,18.7$ at $512^3$: a plateau, not a turnover), while modes at $\lambda\lesssim0.9\,\lambda_{\rm J}$ grow with transverse resolution and are numerical. The converged edge does not migrate, so the $\lambda_{\rm max}\lesssim1\,\lambda_{\rm J}$ upper bound is resolution-stable, and an independent seed reproduces it to $3\%$. Across the eight runs $\lambda_{\rm max}=0.9$--$1.1\,\lambda_{\rm J}$ ($\lambda_{\rm max}/W_{\rm FWHM}=1.5$--$2.0$; Table~\ref{tab:wave}), so $\lambda_{\rm max}/\lambda_{\rm IM92}=0.4$--$0.5$; extending to $f=1.8,2.0$ gives the same scale, covering the HGBS range ($f\approx0.7$--$1.9$). Periodic and reflecting boundaries are indistinguishable (KS $p=0.17$).

\emph{Comparison with linear theory, and the gravity solver.} Figure~\ref{fig:dispersion} overplots the from-scratch--solved linear dispersion relation of the isothermal equilibrium cylinder (Ostriker, $m=0$; reproducing \citealt{Nagasawa1987,Inutsuka1992,Fischera2012} to $<1\%$: peak $\lambda_{\rm max}=22.0\,H$, cutoff $11.1\,H$). In our normalisation $\lambda_{\rm IM92}=22\,H=2.3\,\lambda_{\rm J}$, so the equilibrium spectrum peaks at $2.3\,\lambda_{\rm J}$ and cuts off at $\approx1.16\,\lambda_{\rm J}$, whereas the collapsing filament's measured spectrum rises monotonically through the equilibrium cutoff without turning over in the converged band; the small-scale rise is numerical because it \emph{diverges with resolution}. Self-gravity uses a periodic FFT with reflecting transverse hydrodynamic boundaries; a transverse domain-doubling ($2\to4\,\lambda_{\rm J}$) test and an isolated-boundary re-solution of Poisson's equation both leave $\lambda_{\rm max}$ unchanged (agreeing to $<1\%$ at $\lambda_{\rm max}$), so the selected wavelength is unbiased by the gravity boundary treatment.

\begin{figure}
\includegraphics[width=0.48\textwidth]{figures/fig_dispersion.pdf}
\caption{The from-scratch--solved eigenvalue dispersion relation of the isothermal equilibrium cylinder (\citealt{Nagasawa1987,Inutsuka1992}), validated to $<1\%$ against \citet{Nagasawa1987} and \citet{Fischera2012} (not a fit): it peaks at $\lambda_{\rm IM92}\approx2.3\,\lambda_{\rm J}$ and vanishes at the cutoff $\lambda_{\rm crit}\approx1.2\,\lambda_{\rm J}$, overplotted on the measured growth-rate spectrum ($f=1.1$, three resolutions, normalised at $\lambda=1.14\,\lambda_{\rm J}$). The measured spectrum instead rises monotonically into the grid-scale region (shaded), confirming the small-scale rise is numerical. The measured wavelength ($\approx1.1\,\lambda_{\rm J}$; $\lambda_{\rm max}/W_{\rm FWHM}\approx2$) lies a factor $\approx2.5$--$3$ below the analytic peak ($5$--$6\times$FWHM) and near the cutoff, so we treat it as an order-$\lambda_{\rm J}$, sub-classical value, not a precisely-located maximum.}
\label{fig:dispersion}
\end{figure}

\begin{figure}
\centering
\includegraphics[width=0.99\columnwidth]{figures/fig_eqrecovery.png}
\caption{Equilibrium-recovery control (the primary validation that the sub-classical scale is physical). The identical growth-rate \emph{measurement} pipeline is applied to a \emph{relaxed, force-balanced} pressure-confined Ostriker equilibrium (blue; prepared by explicit drag relaxation, then evolved drag-free) and to a \emph{dynamically contracting} supercritical filament (orange); the two differ in their prepared initial state, not in the measurement, isolating the effect of starting from equilibrium versus contraction. \emph{Left:} onset fragmentation power spectrum; the equilibrium peaks at the classical $\lambda_{\rm m}\approx22H$ ($\approx5.4\,W_{\rm FWHM}$) while the contracting filament peaks sub-classically ($\approx1.4\,W_{\rm FWHM}$, upper limit). \emph{Right:} the dispersion relation $\Gamma(\lambda)$ versus $\lambda/22H$; the equilibrium shows the classical peak-and-decline with its maximum at $22H$, whereas the contracting filament's $\Gamma$ rises monotonically. The pipeline therefore recovers the classical wavelength when the physics is classical, confirming the sub-classical scale of contracting filaments is physical.}
\label{fig:eqrecovery}
\end{figure}

\emph{Nonlinear bead spacing and timescale separation.} Because the growth rate is fitted while the filament contracts, we checked that the linear peak is the scale at which cores actually form. Evolved into the non-linear regime, the near-critical runs develop discrete beads at $\lambda_{\rm bead}\approx1.1$--$1.6\,\lambda_{\rm J}$ (before merging reduces the count), overlapping the growth-rate range $\lambda_{\rm max}=0.89$--$1.14\,\lambda_{\rm J}$ (Table~\ref{tab:wave}) at the order-$\lambda_{\rm J}$ level. The growth is fast (the mode e-folds $\sim3$--$4$ times while $C$ rises from $\approx20$ to $\gtrsim100$, $\tau_{\rm growth}/\tau_{\rm bg}\approx0.3$--$0.8$), so we tested whether the spectrum depends on the fit window: restricting to the early, near-static phase ($\delta_{\rm rms}<0.05$, $N_J\gtrsim5$) reproduces the full-window spectrum to better than $1\%$ at every physical mode with the peak mode unchanged, at $f=1.1$ and $1.5$. The selected wavelength is therefore set at the low-contrast start of the window and is fit-window invariant, corroborated by the nonlinear bead spacing.

\emph{Pipeline benchmark against the analytic Jeans rate.} To verify that the growth-rate machinery returns physical rates, we seeded a near-uniform self-gravitating box ($\rho\simeq\rho_0$) and measured $\Gamma$ per axial mode. For the well-resolved long-wavelength modes the measured rates match the analytic uniform-medium value $\Gamma(k)=\sqrt{4\pi G\rho_0-c_s^2k^2}$ to $1$--$7\%$ ($n{=}1$: $6.16$ vs $6.23$; $n{=}2$: $6.24$ vs $6.08$; $n{=}4$: $5.80$ vs $5.44$), confirming the pipeline. This also explains the apparently large filament rates ($\Gamma\approx18$--$26$): $\Gamma(\rho_0)=\sqrt{4\pi G\rho_0}\approx6.3$, so the filament value corresponds to $\rho/\rho_0\approx(18/6.3)^2\approx8$, i.e.\ $\Gamma$ is the \emph{local} free-fall rate at the elevated central density during the fit window, not an anomalous super-free-fall rate.

\emph{Robustness to line-mass, initial profile and field strength.} The order-$\lambda_{\rm J}$ result is not specific to one configuration. Across $\beta=0.3,0.5,1.0$ (longitudinal, $f=1.1$) $\lambda_{\rm max}=1.00,1.14,0.80\,\lambda_{\rm J}$ --- order $\lambda_{\rm J}$ with only a weak, non-monotonic $\beta$-dependence within the mode-quantisation step ($\Delta\lambda\approx0.13$--$0.15\,\lambda_{\rm J}$), not a resolved trend. Replacing the Plummer-like ($p\simeq2$) profile with an Ostriker ($p=4$) profile gives $\lambda_{\rm max}=0.80\,\lambda_{\rm J}$ at $128^3$ and $256^3$ --- again order $\lambda_{\rm J}$, so the scale is not a profile-index artefact. But this $p=4$ run is only \emph{approximately} in equilibrium (its central contrast relaxes by $\sim10\times$ before beading), so it too fragments from a contracted state; we therefore report the sub-classical scale as a generic outcome of fragmentation from a non-equilibrium state, not as proof that a true equilibrium cylinder cannot fragment classically (which needs the relaxed control below).

\emph{Recovery of the classical limit for a genuine equilibrium --- the primary control.} The critical test is whether the \emph{identical} pipeline recovers the classical $\lambda_{\rm m}=22H\approx5$--$6\,W_{\rm FWHM}$ for a force-balanced equilibrium cylinder. It does. Because the periodic FFT solver enforces the Jeans swindle at the box-mean density, the analytic Ostriker profile is not the solver's equilibrium (initialised directly it breathes by factors of two), so we relaxed a pressure-confined Ostriker profile to hydrostatic balance with an explicit velocity drag, verified by a no-perturbation control (central density drifts only $\sim3\%$, in the $k=0$ mode), and seeded it with the identical 12-mode perturbation ($L_x=16\,\lambda_{\rm J}$, two resolutions). Its $\Gamma(\lambda)$ shows the textbook peak-and-decline at $\lambda_{\rm max}\approx22H$: $\lambda_{\rm max}/W_{\rm FWHM}=5.0$--$5.5$, $\lambda_{\rm max}/22H=1.07$--$1.10$ (converged; $22H/W_{\rm FWHM}=4.96$--$5.03$ matches the analytic Ostriker $5.07$ to $\approx1\%$); the $\approx7$--$10\%$ long peak is a conservative bias. The \emph{same} pipeline on the contracting filaments gives $\lambda_{\rm max}\approx1\,\lambda_{\rm J}\approx1.4\,W_{\rm FWHM}$ ($\approx0.4\times22H$) with a monotonically rising $\Gamma$ (Fig.~\ref{fig:eqrecovery}). The two differ only in their \emph{prepared initial state}, not the drag-free measurement, so the pipeline is demonstrably unbiased and the sub-classical spacing is a physical consequence of contraction; the closed-loop synthetic test independently recovers injected $\lambda/W$ up to $5.7$ to $<0.5\%$. The control validates the \emph{sign} of the deficit; the contracting $\lambda_{\rm max}$ itself remains an upper limit.

\begin{table}
\caption{Growth-rate (dominant-mode) measurements. $\lambda_{\rm max}$ is the peak of $\Gamma(\lambda)$, $\Gamma_{\rm peak}$ its growth rate, $W$ the per-run Gaussian formation FWHM ($\approx0.55$--$0.67\,\lambda_{\rm J}$; all in $\lambda_{\rm J}$). The $f=1.8,2.0$ rows extend the same method into the supercritical regime (single resolution) and give the same sub-classical scale as the bead-count ensemble ($\lambda/\lambda_{\rm J}\approx0.98$), so the two methods agree at the $\sim\lambda_{\rm J}$ level. The discrete axial modes ($\lambda=L_x/n$, $\Delta\lambda\approx0.15\,\lambda_{\rm J}$ near $\lambda\sim1$) give $\lambda_{\rm max}$ a quantisation uncertainty $\pm0.07\,\lambda_{\rm J}$; the final quoted digit is not significant, and the box cannot be lengthened to refine it because $L_x\geq24\,\lambda_{\rm J}$ introduces a seeding artefact.}
\label{tab:wave}
\footnotesize
\setlength{\tabcolsep}{4pt}
\begin{tabular}{cccccc}
\toprule
$f$ & resolution & seed & $\lambda_{\rm max}$ & $\Gamma_{\rm peak}$ & $\lambda_{\rm max}/W$ \\
\midrule
1.1 & $128^3$ & 42  & 1.14 & 17.8 & 1.95 \\
1.1 & $256^3$ & 42  & 1.14 & 18.5 & 1.99 \\
1.1 & $512^3$ & 42  & 1.14 & 18.7 & 2.00 \\
1.2 & $256^3$ & 42  & 1.00 & 19.1 & 1.81 \\
1.2 & $256^3$ & 137 & 0.89 & 20.1 & 1.60 \\
1.5 & $256^3$ & 42  & 0.89 & 21.2 & 1.47 \\
1.5 & $256^3$ & 137 & 0.89 & 22.1 & 1.47 \\
1.5 & $512^3$ & 42  & 0.89 & 21.5 & 1.47 \\
1.8 & $256^3$ & 42  & 0.80 & 24.6 & 1.19 \\
2.0 & $256^3$ & 42  & 0.80 & 26.3 & 1.20 \\
\bottomrule
\end{tabular}
\end{table}

A complementary end-to-end grid of 21 ambient-confined runs ($f=1.5$--$3.0$, $\beta=0.5$--$2.0$, two seeds; campaigns~11--13), forward-modelled through beam convolution and Plummer fitting, gives a consistent picture: $\lambda/\lambda_{\rm J}(\rho_0)\approx0.9$--$1.2$, with boundary condition, profile and seed each moving the value by $\lesssim40\%$ (medians $1.2$ reflecting, $0.8$ periodic, $1.4$ user; $1.4$ Gaussian, $0.9$ Ostriker) and a factor-$1.6$ contrast variation shifting it by only $\sim0.3$, so the wavelength is a property of the filament, not the confinement. All per-run measurements are released with the public archive.

\emph{Numerical validation (summary).} The fragmentation \emph{classification} is resolution-independent ($t_{\rm frag}(256^3)/t_{\rm frag}(128^3)=0.915\pm0.012$; uniform-density ICs give identical classifications). The Jeans length stays resolved where the spacing is measured for the headline longitudinal runs ($N_J\approx5$--$7$; Truelove $N_J\gtrsim4$ satisfied) and the ideal-perpendicular ladder reaches $N_J\approx40$--$170$; only the perpendicular \emph{ambipolar} runs remain unconverged ($N_J\approx1$--$2$), so their short wavelengths are treated qualitatively. Full convergence tables and the complete 2251-run campaign list are in Appendix~\ref{sec:validation} and Table~\ref{tab:campaigns}.


% ===== Reintegrated supporting appendices =====
\section{Simulation Validation Details}
\label{sec:validation}

To ensure robustness against numerical effects we performed three validation tests (83 simulations). For resolution, six parameter points run at both $128^3$ and $256^3$ with matched problem-generator settings, initial conditions, seeds and boundaries all fragment at both resolutions; the mean timescale ratio is
\begin{equation}
\frac{t_{\rm frag}(256^3)}{t_{\rm frag}(128^3)} = 0.915 \pm 0.012,
\end{equation}
an $8.5\%\pm1.2\%$ offset (from the ratio $0.915\pm0.012$), all points within a $\pm11\%$ band (Fig.~\ref{fig:resolution}), so the fragmentation classification is resolution-independent. Replacing the Gaussian-profile initial condition with uniform density yields identical classifications at all 10 tested points.

A fragmentation-wavelength measurement is trustworthy only if the local Jeans length remains resolved ($\gtrsim4$ cells; Truelove et al.\ 1997) up to the measurement epoch, because an under-resolved Jeans length spuriously enhances small-scale power --- exactly the grid-scale, resolution-divergent rise we identify as numerical (Fig.~\ref{fig:growthrate}a). We tracked the minimum $N_J$ along each collapse (Table~\ref{tab:truelove}). For the \emph{longitudinal} runs (near-critical calibration and supercritical ensemble) the inter-bead wavelength ($\lambda\approx\lambda_{\rm J}$, resolved by $\sim32$ cells) is set at moderate contrast ($C\approx20$--$40$) where $N_J\approx5$--$7$; $N_J$ drops below 4 only at $C>170$, after the spacing is set. For the perpendicular \emph{ambipolar} runs, beading is measured only at $C\approx250$--$560$ where $N_J\approx1$--$2$ (Truelove violated), so the short perpendicular values are resolution-limited upper limits. This is why we base the quoted wavelength on the growth-rate spectrum (established at $N_J\gtrsim5$, resolution-converged) rather than late-time peak counts (set by sub-fragmentation at $N_J<4$).

\begin{table}
\caption{Jeans (Truelove) resolution: minimum cells per local Jeans length $N_J$ at the epoch the spacing is measured, and the density contrast $C$ at that epoch.}
\label{tab:truelove}
\footnotesize
\setlength{\tabcolsep}{3pt}
\begin{tabular}{@{}p{3.0cm}p{1.9cm}p{2.5cm}@{}}
\toprule
Configuration & $C$ at meas. & min $N_J$ \\
\midrule
Near-critical longitudinal (calibration) & $20$--$40$ & $5$--$7$ (OK) \\
Supercritical longitudinal ensemble & early beading, moderate $C$ & $\gtrsim4$ (OK) \\
Ideal perp., converged ladder ($512^3$; \S\ref{sec:optional_tests}) & $\approx3$ / $\approx30$ & $170$ / $40$ (OK; monolithic) \\
Ideal perp., short-$\lambda$ artefact onset & $>10^2$ (grid-dep.) & $\lesssim1$ (violated) \\
Perp.\ ambipolar (tension side) & $250$--$560$ & $1$--$2$ (violated) \\
\bottomrule
\end{tabular}
\end{table}

\begin{figure}
\includegraphics[width=0.48\textwidth]{figures/figR1_resolution_scatter.pdf}
\caption{Resolution convergence with matched problem generator: $t_{\rm frag}(256^3)$ versus $t_{\rm frag}(128^3)$ for six parameter points, all within a $\pm11\%$ band (mean ratio $0.915\pm0.012$).}
\label{fig:resolution}
\end{figure}

\subsubsection{Equation of State Sensitivity}

Testing mildly sub-isothermal equations of state with $\gamma = 0.9$ and $0.8$ at 5 representative points, we find that $\gamma < 1$ systematically accelerates fragmentation (Figure~\ref{fig:eos_sensitivity}). At $\gamma = 0.8$, all 5 test points fragmented at $t_{\rm frag} \approx 0.66$--$0.68\,t_{\rm J}$, approximately half the isothermal fragmentation timescale for the same $(f, \beta, \mathcal{M})$; at $\gamma = 0.9$ fragmentation is delayed (1 of 5 fragments by the integration limit). These are relative comparisons at matched $t_{\rm J}$: the point is that $\gamma<1$ shortens $t_{\rm frag}$ monotonically, consistent with the isothermal runs of the main text (which fragment at $t_{\rm frag}\approx1.1$--$1.2\,t_{\rm J}$); we draw no stability inference from the finite integration window.

The opposite limit confirms the mechanism: an \emph{adiabatic} equation of state ($\gamma=5/3$; campaign~8, 30 runs) suppresses fragmentation entirely (no fragmentation to $t>30$--$40\,t_{\rm J}$ at any tested $f$), because compressive heating restores Jeans support. For $\gamma<1$ the pressure response to compression weakens ($P\propto\rho^{\gamma}$ rises more slowly than isothermal), reducing Jeans support and hastening fragmentation. Far-IR-cooled cloud filaments typically have $\gamma\approx0.7$--$0.95$, so the isothermal predictions are conservative: real filaments are if anything more fragmentation-prone. The $\gamma<1$ and $\gamma=5/3$ limits bracket the isothermal case and identify the equation of state as what would set a small-scale fragmentation cutoff --- the point relevant to the opacity limit (Section~\ref{sec:limitations}).

\begin{figure}
\includegraphics[width=0.48\textwidth]{figures/fig6_phase3_eos_sensitivity.pdf}
\caption{Equation of state sensitivity test. Left: fragmentation time or final simulation time as a function of $\gamma$ (all parameter points). Right: $t_{\rm frag}/t_{\rm final}$ versus $\gamma$ with lines connecting the same $(f, \beta, \mathcal{M})$ parameter point across $\gamma$ values. Sub-isothermal EOS ($\gamma < 1$) systematically reduces $t_{\rm frag}$.}
\label{fig:eos_sensitivity}
\end{figure}

\section{Complete simulation campaign list}
\label{sec:campaigns}

Table~\ref{tab:campaigns} lists every simulation campaign cited in the manuscript. Rows 5--8 are the components of what was previously reported as a single Field-Geometry campaign (314 runs); they are counted individually here and the 314 aggregate is \emph{not} added separately.

\begin{table*}
\caption{Complete list of simulation campaigns underlying the paper. The ``Location'' column points, via cross-reference, to the section or appendix where each campaign is discussed in this consolidated manuscript. Campaign~9 aggregates the turbulent-support $f_{\rm eff}$ grid (90 runs), the physical-vs-synthetic turbulence comparison (54 runs), and the perpendicular-$\beta$ and critical-transition grids. ``FG comp.'' denotes a component of the former Field-Geometry campaign (rows 5--8).}
\label{tab:campaigns}
\footnotesize
\setlength{\tabcolsep}{4pt}
\renewcommand{\arraystretch}{1.2}
\begin{tabular}{@{}c p{3.4cm} l p{4.0cm} r l@{}}
\toprule
\# & Campaign & Location & Coverage $(f,\beta,\mathcal{M},\theta,\eta_{\rm AD})$ & Runs & Indep.? \\
\midrule
1  & Base three-regime grid & App.~\ref{sec:supplementary_campaigns_app} & $f\approx1$; $\beta\,0.05$--$5$; $\mathcal{M}\,0.5$--$5$; $\theta=0$; ideal & 208 & yes \\
2  & Transition grid / DTC & App.~\ref{sec:supplementary_campaigns_app} & $f\,1.4$--$2.2$; $\beta\,0.3$--$1.3$; $\mathcal{M}\,1$--$5$; $\theta=0$ & 540 & yes \\
3  & Supercritical grid & \S\ref{sec:modecompetition} & $f\,1.1$--$3.0$; $\beta\,0.3$--$5.0$; $\mathcal{M}\,0.5$--$3$; periodic & 654 & yes \\
4  & Validation & App.~\ref{sec:validation} & resolution / IC / EOS convergence & 83 & yes \\
5  & Near-critical & \S\ref{sec:modecompetition} & $f\,1.0$--$1.2$; $\beta\,0.3$--$1.0$; $\mathcal{M}\,1$--$2$; $\theta=0$ & 80 & FG comp. \\
6  & Perpendicular & \S\ref{sec:modecompetition} & $\theta=90^\circ$; $f\,1.2$--$1.5$; $\beta\,0.3$--$2.0$ & 96 & FG comp. \\
7  & Oblique calibration & App.~\ref{sec:supplementary_campaigns_app} & $\theta=30/45/60^\circ$; $\beta\,0.5$--$2.0$ & 108 & FG comp. \\
8  & Adiabatic EOS & App.~\ref{sec:validation} & $\gamma=5/3$; $f\,1.0$--$1.2$ & 30 & FG comp. \\
9  & Turbulence + perp-$\beta$ + crit.-transition & App.~\ref{sec:supplementary_campaigns} & $\delta v/c_s\,10^{-4}$--$1$; $\beta\,0.3$--$2$ & 289 & yes (agg.) \\
10 & Targeted DTC re-runs & App.~\ref{sec:supplementary_campaigns_app} & $\beta=0.3,\mathcal{M}=1$; extended timeout & 25 & yes \\
11 & Direct supercritical ensemble & \S\ref{sec:modecompetition} & $f\,1.5$--$2.0$; $\beta\,0.5$--$2.0$; $\theta=0$; ambient & 39 & yes \\
12 & Reconciliation suite & \S\ref{sec:modecompetition} & $f\times\beta\times2$ seeds; $\theta=0$; ambient & 12 & yes \\
13 & Ideal-MHD perp.\ $d=1.0$ grid & \S\ref{sec:modecompetition} & $\theta=90^\circ$; $\beta\,0.5$--$2.0$; ambient & 18 & yes \\
14 & Ambipolar non-ideal survey & \S\ref{sec:modecompetition} & $\theta=90^\circ$; $\eta_{\rm AD}\,0.001$--$0.1$; $\beta\,0.5$--$2.0$ & 36 & yes \\
15 & Growth-rate + equil.\ control + transverse/higher-$f$ tests & \S\ref{sec:validation_endtoend}, App.~\ref{sec:val_appendix} & $f\,0.7$--$2.0$; $128^3$--$512^3$; Gaussian/Ostriker; transverse $2$--$4\,\lambda_{\rm J}$ & 25 & yes \\
16 & Axial-resolution stability & \S\ref{sec:optional_tests} & $f=1.1$; $\theta=0$; $1024\times128^2$ rerun & 1 & yes \\
17 & Ideal-perp.\ refinement ladder & \S\ref{sec:optional_tests} & $\theta=90^\circ$; $f=1.1$; $128^3$--$512^3$ & 3 & yes \\
18 & Divergence-free helical grid & \S\ref{sec:optional_tests} & $f=1.1$; $B_\phi/B_z=0$--$2$; $512\times128^2$ & 4 & yes \\
\midrule
   & \textbf{Athena++ MHD total} & & & $\boldsymbol{2251}$ & \\
\bottomrule
\end{tabular}
\end{table*}

\section{Supplementary simulation-campaign results}
\label{sec:supplementary_campaigns_app}

The campaigns in this appendix are exploratory and are \emph{not} used quantitatively in the main conclusions; the headline wavelength rests solely on the near-critical longitudinal growth-rate ladder of Section~\ref{sec:validation_endtoend}. For completeness, the \emph{oblique}-field campaign (108 runs, $\theta=30/45/60^\circ$) shows the collapse time decreasing monotonically with field angle ($t_{\rm frag}=0.60,0.51,0.45\,t_{\rm J}$), reflecting accelerating \emph{radial} collapse as axial support weakens rather than a shortening of the axial mode; with three angles and no per-angle error this is empirical consistency, not a calibrated $\lambda_{MJ}(\theta,\beta)$ relation, so we draw no oblique $\lambda/W$ conclusion.

\subsection{The radial-collapse/fragmentation timescale competition}
\label{sec:regimes}

The fate of an isothermal, self-gravitating filament is governed by the competition between radial collapse and longitudinal fragmentation \citep{Inutsuka1997,Pon2012}. Near-critical filaments ($f\lesssim1.2$) retain enough support to evolve quasi-statically, so the longitudinal instability grows to non-linearity and $\lambda/W$ is measurable directly. Supercritical filaments ($f\gtrsim1.5$) fragment longitudinally when confined by an ambient medium (Section~\ref{sec:modecompetition}), at a converged \emph{upper bound} $\lesssim0.5$ times the classical value; without confinement the Gaussian skirt feeds an accretion channel that suppresses fragmentation. Extended-integration tests confirm the periodic grid's early termination fires before the instability becomes non-linear ($\sigma_{\rm lon}$ onset at $t=0.42$--$0.71\,t_{\rm J}$ versus a Courant kill at $t\approx0.25$--$0.34\,t_{\rm J}$), independent of resolution.

In the near-critical regime ($f\approx1.0$--$1.2$) filaments bead at $t_{\rm frag}\approx1.1$--$1.2\,t_{\rm J}$; runs at the exact Jeans-critical value ($f=1.00$) also bead, at $t_{\rm frag}=1.62\pm0.03\,t_{\rm J}$ (later than the more strongly seeded $f=1.0$--$1.2$ grid), so there is no stability threshold at the critical line mass and $\lambda/W$ is directly measurable here.

\textbf{Supercritical regime ($f\geq1.5$): radial collapse dominates.} In the periodic supercritical grid ($f=1.5$--$3.0$) we find no longitudinal structure, only pure radial collapse (central density rising two orders of magnitude while the longitudinal variance stays $<0.02\%$), with runaway ($\Delta t<10^{-8}\,t_{\rm J}$) before beading develops. The radial-collapse time ($t_{\rm coll}\approx0.25$--$0.34\,t_{\rm J}$) is $3$--$5\times$ shorter than the near-critical fragmentation time ($1.1$--$1.6\,t_{\rm J}$), so out-of-equilibrium configurations reach runaway before longitudinal perturbations grow. The measurable wavelengths therefore come from the near-critical quasi-static regime and from ambient-confined supercritical filaments that fragment before radial collapse (Section~\ref{sec:modecompetition}).

\subsection{Three fragmentation regimes}

The baseline grid ($f\approx1$) reveals three fragmentation regimes in the $(\beta,\mathcal{M})$ plane, characterised by the end-of-run density contrast $C=\rho_{\rm max}/\rho_0$ (Fig.~\ref{fig:regime}): (i) \emph{strong-field, magnetically supported} ($\beta\lesssim0.15$; $C\approx1$); (ii) \emph{magnetically regulated} ($0.2\lesssim\beta\lesssim2$; $C=3$--$11$); and (iii) \emph{thermally dominated} ($\beta\gtrsim3$; $C=13$--$22$). The dispersion relation \citep{Nakamura1993} predicts near-complete suppression for $\beta\lesssim0.15$ and $<10\%$ wavelength modification for $\beta\gtrsim2$--$3$, consistent with the empirical boundaries (which bracket the sampled grid points at the one-grid-step level, so the thresholds carry $\sim$one-grid-step uncertainty). Observed dense-filament conditions ($\mathcal{M}\approx1$--$3$, $\beta\approx0.3$--$2$; \citealt{Arzoumanian2019,Planck2016,Pattle2018}) place most HGBS filaments in the magnetically regulated regime.

\begin{figure}
\includegraphics[width=\columnwidth]{figures/fig2_regime_corrected.pdf}
\caption{Three-regime fragmentation framework from the Athena++ base grid (near-critical, $\theta=0^\circ$, isothermal; Section~\ref{sec:regimes}). Colour shows the end-of-run density contrast $C=\rho_{\rm max}/\rho_0$. The regimes, set by plasma $\beta$: (I) \emph{strong-field, magnetically supported} ($\beta\lesssim0.15$; $C\approx1$); (II) \emph{magnetically regulated} ($0.2\lesssim\beta\lesssim2$; $C=3$--$11$); (III) \emph{thermally dominated} ($\beta\gtrsim3$; $C=13$--$22$). Typical HGBS conditions ($\mathcal{M}\approx1$--$3$, $\beta\approx0.3$--$2$) lie in the magnetically regulated regime. This $\beta$ criterion is distinct from the magneto-static mass-to-flux criticality (Section~\ref{sec:regimes}).}
\label{fig:regime}
\end{figure}

\subsection{Turbulence and critical-transition dependence}
\label{sec:supplementary_campaigns}

Additional grids map the turbulence, field-geometry and critical-transition dependence of $\lambda/W$.

\subsubsection{Turbulence versus $\lambda/W$}

Comparing physical turbulence ($\delta v/c_s=1.0$, Kolmogorov) with synthetic perturbations ($10^{-4}$) at longitudinal field ($f=1.0$--$1.2$, $\beta=0.5,1,2$), the fragmentation \textit{wavelength} is unchanged ($<5\%$; $\lambda/W=3.44\pm0.76$ vs $3.30\pm0.85$; Fig.~\ref{fig:C5_turbulence}), with a $\beta$-dependence only in the raw late-time bead-count ratio ($4.31\pm0.60$ at $\beta=0.5$ to $2.81\pm0.13$ at $\beta=2.0$), which the growth-rate analysis (Section~\ref{sec:validation_endtoend}) shows is a late-epoch effect, not a shift in $\lambda_{\rm max}$. Turbulence accelerates collapse by $3.2\times$ ($t_{\rm frag}=0.33\pm0.06$ vs $1.06\pm0.16\,t_{\rm J}$). As a single decaying realisation this does not resolve whether \emph{driven} transonic turbulence is a first-order parameter; we conclude only that the spacing is insensitive to the seed amplitude in this limited test.

\begin{figure}
\includegraphics[width=\columnwidth]{figures/fig_C5_turbulence_lambdaW.pdf}
\caption{turbulence comparison: physical vs synthetic turbulence give statistically indistinguishable $\lambda/W_{\rm core}$ ($\Delta\mu<5\%$, KS $p=0.43$). The $\beta$-trend in the raw late-time bead-count ratio ($\sim$4.3 at $\beta=0.5$ to $\sim$2.8 at $\beta=2.0$) is preserved across the two turbulence prescriptions, but is a late-epoch bead-count effect rather than a shift in the underlying growth-rate $\lambda_{\rm max}$, whose $\beta$-variation is weak (Section~\ref{sec:validation_endtoend}).}
\label{fig:C5_turbulence}
\end{figure}

\subsubsection{Perpendicular-field spacing}

For perpendicular fields ($\theta=90^\circ$, $f=1.2$--$1.5$, $\beta=0.3$--$2.0$; extended $16\lambda_{\rm J}$ periodic axial domain): strong perpendicular B ($\beta\leq0.5$) gives pure radial collapse with no axial fragmentation, while for $\beta\geq1.0$ axial beading occurs at $\lambda/W_{\rm core}=1.25\pm0.09$ ($N=27$), i.e.\ $\lambda/W_{\rm fil}\approx0.4$ under the corrected conversion (Eq.~\ref{eq:t1_apply}), a factor $2$--$3$ \textit{shorter} than the longitudinal-field value at the same $\beta$. Field geometry, not $\beta$, dominates, and the perpendicular spacing lies far below the observed NN value (fixed-width, periodic-corrected, $\lambda/W\approx1.9$; raw $2.1$).

The two perpendicular campaigns give different $\beta$-trends: this periodic-domain grid beads only for $\beta\geq1.0$, whereas the ambient-confined ideal-MHD grid of Section~\ref{sec:modecompetition} beads at $\beta=0.5$ and $2.0$ with spindle collapse near $\beta\approx1$. They differ in transverse confinement and are both numerically unresolved ($N_J\approx1$--$2$; Appendix~\ref{sec:validation}), so we draw no physical $\beta$-trend from either; a matched-resolution, matched-confinement grid is needed to settle it (Section~\ref{sec:limitations}). The robust statement is only that the perpendicular beading wavelength, where it occurs, is short ($\lambda/W_{\rm fil}\lesssim0.6$), well below the observed spacing.

\subsubsection{Critical-transition spacing}

Mapping $\lambda/W$ across $f=0.9$--$1.3$ at five $\beta$ values (Fig.~\ref{fig:C7_critical}) shows no discontinuity at $f=1.0$ (variation $<5\%$) and an apparent $\beta$-dependence in the \emph{late-time bead count}: $\lambda/W_{\rm core}=4.74,4.38,3.19,2.86,2.80$ at $\beta=0.3,0.5,1.0,1.5,2.0$. This is larger than the primary growth-rate measurement ($\lambda_{\rm max}=1.00,1.14,0.80\,\lambda_{\rm J}$ at $\beta=0.3,0.5,1.0$, a weak mode-quantised trend), so we treat the wavelength's $\beta$-dependence as weak (the stronger bead-count trend being a late-epoch, resolution-sensitive artefact), with the highest-$\beta$ points ($\lambda/\lambda_{\rm J}\approx0.84$) consistent with the supercritical ensemble. \textbf{Continuity across $f\approx1$ does not validate extrapolation to $f\geq1.5$}, where the mode changes to radial collapse; the near-critical $\lambda/W$ values are a hypothesis for supercritical filaments, not a measurement.

\begin{figure}
\includegraphics[width=\columnwidth]{figures/fig9_nearcritical_corrected.pdf}
\caption{Near-critical runs: $\lambda/W_{\rm core}$ across $f=0.9$--$1.3$ for five $\beta$ values, smooth with no discontinuity at $f=1.0$ (left); mean per $\beta$ (right). The convergence-independent quantity is $\lambda/W_{\rm core}$; the observable follows from $\lambda/W_{\rm fil}\approx0.27\,(\lambda/W_{\rm core})$ (Eq.~\ref{eq:t1_apply}), so the highest-$\beta$ points correspond to $\lambda/W_{\rm fil}\approx0.75$--$0.8$. This near-critical calibration lies below the observed spacing under any convention (Table~\ref{tab:headline}); the primary simulated value is the growth-rate wavelength (an upper limit; Section~\ref{sec:validation_endtoend}), not this ratio.}
\label{fig:C7_critical}
\end{figure}

\begin{table}
\caption{Effective line-mass fraction with turbulent support.}
\label{tab:feff}
\footnotesize
\setlength{\tabcolsep}{4pt}
\begin{tabular}{p{3.0cm}cc}
\toprule
Quantity & Value & Source \\
\midrule
Nominal $f$ (thermal) & $1.5 \pm 0.3$ & HGBS line-mass \\
Turbulent support $f_{\rm eff}/f$ & $0.82 \pm 0.11$ & TURB (90 sims) \\
Effective $f_{\rm eff}$ & $1.23 \pm 0.28$ & product \\
HGBS ensemble range & $0.7$--$1.9$ & propagated \\
\bottomrule
\end{tabular}
\end{table}

\subsubsection{Synthesis}

Table~\ref{tab:cross_campaign} summarises the simulation $\lambda/W$ measurements and their role in the comparison with observation.

\begin{table*}
\caption{Simulation fragmentation scale by configuration, in the two convention-independent forms: the raw internal ratio $\lambda/W_{\rm core}$ (normalised by the initial half-width) and the background-density Jeans-length ratio $\lambda/\lambda_{\rm J}=0.3\,(\lambda/W_{\rm core})$. We do \emph{not} tabulate the observable $\lambda/W_{\rm fil}$ here, since it depends on the width convention; the comparison with observation uses the growth-rate wavelength expressed in \emph{matched} beam-convolved Plummer units, $\lambda/W\approx1.0$--$1.3$ ($\approx0.35$--$0.5$ times the classical value, a factor $\approx2$--$3$; Section~\ref{sec:validation_endtoend}), which is set against the observed $\lambda/W\approx0.4$ (FilFinder) to $1.9$ (fixed-width DisPerSE).}
\label{tab:cross_campaign}
\footnotesize
\setlength{\tabcolsep}{4pt}
\begin{tabular}{lccl}
\toprule
Configuration & $\lambda/W_{\rm core}$ & $\lambda/\lambda_{\rm J}$ & Status \\
\midrule
\multicolumn{4}{l}{\emph{Longitudinal field (quantitative)}} \\
Long.\ near-critical ($\beta{=}2.0$, calibration) & $2.80{\pm}0.14$ & $0.84$ & Calibration only \\
Long.\ supercritical ensemble (\S\ref{sec:modecompetition}) & $3.27$ & $0.98$ & Near one Jeans length; see growth-rate result \\
\midrule
\multicolumn{4}{l}{\emph{Perpendicular field --- PRELIMINARY, numerically unconverged ($N_J{\approx}1$--$2$); not for quantitative use}} \\
Perp.\ near-critical ($\beta{\geq}1$, ideal-MHD) & $[1.25]$ & $[0.38]$ & unconverged \\
Perp.\ ambipolar ($\eta_{\rm AD}{=}0.001$--$0.1$) & $[{\sim}2.2]$ & $[0.66]$ & unconverged \\
\bottomrule
\end{tabular}
\\[2pt]
Only the longitudinal rows are used quantitatively, and even there the reference value is the convergently-bounded growth-rate wavelength (an upper limit; Section~\ref{sec:validation_endtoend}, Table~\ref{tab:wave}), not these raw ratios. The perpendicular values are bracketed to signal that they are numerically unconverged ($N_J\approx1$--$2$, below the Truelove criterion; wavelengths shift by factors $2$--$5$ between resolutions) and must not be read as measurements.
\end{table*}

Except for the 54-run physical-versus-synthetic turbulence \emph{comparison} (distinct from the 90-run turbulent-\emph{support} $f_{\rm eff}$ grid of Table~\ref{tab:feff}, both in the 289-run campaign~9), the grid uses near-laminar seeding ($\delta v=\mathcal{M}c_s\times10^{-4}$) rather than driven ISM turbulence; that comparison shows the wavelength is unchanged ($<5\%$) at driven amplitudes, though the timescales shorten by $3$--$5\times$. In convergence-independent units the two first-order parameters are field geometry (perpendicular $\lambda/W_{\rm core}\approx1.25$ versus longitudinal $2.8$--$4.4$) and, for longitudinal fields, plasma $\beta$; the observable comparison uses the growth-rate wavelength (Section~\ref{sec:validation_endtoend}), not these raw ratios.

\bibliography{references_complete}

\end{document}
